\documentclass[11pt,letterpaper]{article}
\usepackage[T1]{fontenc}
\usepackage[utf8]{inputenc}
\usepackage{lmodern}
\usepackage[margin=1in,headheight=15pt]{geometry}
\usepackage{amsmath,amssymb,amsthm,mathtools,mathrsfs}
\usepackage{microtype,booktabs,longtable,array,enumitem}
\usepackage[dvipsnames]{xcolor}
\usepackage{fancyhdr,aliascnt}
\usepackage[colorlinks=true,linkcolor=MidnightBlue,citecolor=MidnightBlue,urlcolor=MidnightBlue]{hyperref}
\usepackage[nameinlink,noabbrev]{cleveref}
\hypersetup{pdftitle={Jordan Separation, Cauchy Theory, and Stokes Calculus on the Surcomplex Plane},pdfauthor={Research exposition},pdfsubject={Surcomplex topology, Hahn integration, semialgebraic winding numbers, and microscopic residues}}
\pagestyle{fancy}\fancyhf{}
\fancyhead[L]{\small Surcomplex contours and Stokes calculus}
\fancyhead[R]{\small Topology, chains, and residues}
\fancyfoot[C]{\thepage}
\setlength{\parskip}{0.3em}
\setlength{\parindent}{1.3em}
\setlist{nosep,leftmargin=2em}
\allowdisplaybreaks[2]
\newtheorem{theorem}{Theorem}[section]
\newaliascnt{lemma}{theorem}\newtheorem{lemma}[lemma]{Lemma}\aliascntresetthe{lemma}
\newaliascnt{proposition}{theorem}\newtheorem{proposition}[proposition]{Proposition}\aliascntresetthe{proposition}
\newaliascnt{corollary}{theorem}\newtheorem{corollary}[corollary]{Corollary}\aliascntresetthe{corollary}
\theoremstyle{definition}
\newaliascnt{definition}{theorem}\newtheorem{definition}[definition]{Definition}\aliascntresetthe{definition}
\newaliascnt{example}{theorem}\newtheorem{example}[example]{Example}\aliascntresetthe{example}
\theoremstyle{remark}
\newaliascnt{remark}{theorem}\newtheorem{remark}[remark]{Remark}\aliascntresetthe{remark}
\crefname{theorem}{theorem}{theorems}\crefname{lemma}{lemma}{lemmas}
\crefname{proposition}{proposition}{propositions}\crefname{corollary}{corollary}{corollaries}
\crefname{definition}{definition}{definitions}\crefname{example}{example}{examples}
\crefname{remark}{remark}{remarks}
\newcommand{\No}{\mathbf{No}}\newcommand{\SC}{\mathbf{SC}}
\newcommand{\R}{\mathbb R}\newcommand{\C}{\mathbb C}\newcommand{\Q}{\mathbb Q}
\newcommand{\N}{\mathbb N}\newcommand{\Z}{\mathbb Z}
\newcommand{\HH}{\mathcal H}\newcommand{\A}{\mathscr A}\newcommand{\OO}{\mathcal O}
\newcommand{\m}{\mathfrak m}\newcommand{\st}{\operatorname{st}}
\newcommand{\supp}{\operatorname{supp}}\newcommand{\Res}{\operatorname{Res}}
\newcommand{\res}{\operatorname{res}}\newcommand{\Wind}{\operatorname{Wind}}
\newcommand{\ord}{\operatorname{ord}}\newcommand{\Tr}{\operatorname{Tr}}
\newcommand{\id}{\operatorname{id}}\newcommand{\vH}{v_{\!H}}
\newcommand{\HInt}{\mathop{\int^{\!H}}\nolimits}
\newcommand{\HCoint}{\mathop{\oint^{\!H}}\nolimits}
\newcommand{\sh}{^{\sharp}}\newcommand{\dd}{\,d}
\newcommand{\abs}[1]{\left|#1\right|}
\newcommand{\set}[1]{\left\{#1\right\}}
\newcolumntype{L}[1]{>{\raggedright\arraybackslash}p{#1}}
\numberwithin{equation}{section}
\begin{document}
\begin{titlepage}
\centering
\vspace*{0.7cm}
{\Large\scshape A research synthesis and continuation\par}
\vspace{0.8cm}
{\Huge\bfseries Jordan Separation,\par Cauchy Theory,\par and Stokes Calculus\par}
\vspace{0.45cm}
{\LARGE on the Surcomplex Plane\par}
\vspace{0.7cm}
{\large Topologies, admissible chains, multiscale contours,\par and radius-free residues\par}
\vspace{0.6cm}
{\large September 21, 2026\par}
\vspace{0.65cm}
\begin{minipage}{0.94\textwidth}
\begin{abstract}
A counterpart of a complex-analytic theorem over $\SC=\No[i]$ must specify both a topology and an integration mechanism. The full fine topology has no nonconstant continuous real-parameter paths; a fixed Hahn field has a different, nondiscrete but totally disconnected internal topology. Neither supports the literal ordinary Jordan-curve picture. We distinguish two useful replacements: Jordan separation in the coarse standard-part topology, and the established semialgebraic or definable Jordan theorem over real closed fields.

For integration we develop a Hahn-valued de Rham complex on ordinary manifolds, prove generalized Stokes and a cohomology comparison, and extend the calculus to support-controlled infinitesimal deformations of chains. Cauchy's theorem becomes a closed-form period statement, not a consequence of fine connectedness. An exact endpoint formula describes deformed contour integrals. A multiscale rescaling theorem then turns radius-free Hahn analytic germs into polynomial coefficient families. It realizes one-variable and multidimensional perturbation residues by microscopic contour functionals; a determinant transformation handles arbitrary isolated complete intersections.

We also construct a rational contour pairing from semialgebraic winding numbers and compare it with Hahn integration where both are defined. Existing semialgebraic integration, Berkovich--tropical Stokes formulas, surreal antiderivative theories, and roots-of-unity quadrature are assessed separately. Explicit examples show infinitesimal area, cancellation of infinite residues, and failure of naive finite-mesh or valuation-limit definitions.
\end{abstract}
\end{minipage}
\vfill
\begin{minipage}{0.94\textwidth}\small
\textbf{Status.} Established results, statements reported in the two supplied manuscripts, and deductions proved here are distinguished throughout. This is not a claim of a unique foundation for surcomplex analysis, an exhaustive priority determination, or a formally verified proof. All sums and mathematical workspaces used in the constructions are set-sized.
\end{minipage}
\end{titlepage}
\begingroup
\setlength{\parskip}{0pt}
\tableofcontents
\endgroup
\newpage

\section{The question and the principal conclusions}\label{sec:introduction}

The three classical theorems in the question concern different structures. The Jordan curve theorem describes separation by a topological embedding of a circle. Cauchy's integral theorem describes the vanishing of periods of holomorphic differential forms under appropriate homological hypotheses. Generalized Stokes relates a differential on forms to a boundary on chains. Even over $\C$, an arbitrary Jordan curve need not be rectifiable, and a contour need not be simple for Cauchy's theorem to apply. Thus one should not try to extend the three theorems by a single replacement of $\C$ with a larger field.

The useful conclusions are these. A genuine Jordan analogue is available in \emph{semialgebraic topology} over every real closed Hahn workspace. A second, intentionally coarse Jordan statement is available after passing through standard parts. An exact $\SC$-valued Stokes calculus is available on \emph{Hahn-coherent chains}, including infinitesimal deformations and microscopic affine charts. Finally, rational differentials admit a contour pairing based on algebraic winding numbers even when an ordinary leading contour meets the reduction of a pole. These answers coexist, but they are not interchangeable.

\subsection{What the supplied manuscripts actually provide}
We refer to \emph{Surcomplex Analysis: Infinitesimal Calculus, Hahn-Coherent Holomorphy, and Contour Theory} as the \emph{analysis manuscript} \cite{Analysis}. It distinguishes fine-local power-series calculus from the algebra
\[
 \HH(U)=\OO(U)((t^\Gamma))
\]
of well-ordered Hahn families of ordinary holomorphic functions on a common ordinary domain $U$. Its section ``Contour functionals, Cauchy theory, and primitives'' defines integration coefficientwise on ordinary paths and proves Cauchy, Morera, primitive, and integral-estimate theorems. Its later sections prove moving-pole residue formulas and permit coherent affine changes of scale. These are existing results within the supplied framework, not results first introduced in the present article.

The second source, \emph{Infinitesimal Analytic Geometry over the Surcomplex Numbers}, is the \emph{geometry manuscript} \cite{Geometry}. It instead uses
\[
 \A_n=\C\{z_1,\ldots,z_n\}((t^\Gamma)),
\]
where each coefficient is a convergent complex germ but the coefficients need not have a common radius. For isolated complete intersections it defines a perturbation residue by a strongly summable series of ordinary local residues. It explicitly does \emph{not} define that functional by a surcomplex contour integral. This gap between a well-defined residue and a representing contour is one principal topic below.

Neither manuscript establishes a Jordan theorem for arbitrary fine-topological curves or generalized Stokes for arbitrary fine-smooth functions and arbitrary surcomplex domains. We do not silently attribute either assertion to them. We retain their distinction between ordinary coefficient regularity, strong summability, and fine differentiation.

\subsection{Established external counterparts}
There are relevant results beyond the supplied sources. Berarducci and Otero develop definable intersection numbers, winding numbers, and Jordan--Brouwer separation over o-minimal expansions of real closed fields \cite{BOintersection}. Eisermann gives an algebraic winding-number construction using Cauchy indices and Sturm chains over real closed fields \cite{Eisermann}. Ohmoto and Shiota give $C^1$ semialgebraic triangulations and a Stokes formula; their discussion of non-Archimedean integration explicitly allows a larger saturated target field \cite{OS}. Chambert-Loir and Ducros construct a different Stokes theory for real forms on Berkovich spaces \cite{CLD}. Costin and Ehrlich study surreal extension and integration of distinguished classes of real functions \cite{CE}.

These works answer substantial parts of the question in their own categories. A theorem about real-valued tropical forms is not already a theorem about $K$-valued holomorphic contour integrals. A definably connected interval is not an ordinarily connected interval. An integral valued in a saturated extension is not automatically a canonical element of a previously fixed Hahn field.

\subsection{Deductions developed here}
The main constructions proved in this article are the Hahn de Rham and deformed-chain calculus in \cref{sec:stokes,sec:deformed}; the exact contour-displacement formula in \cref{sec:cauchy}; and the microscopic realization theorems in \cref{sec:radiusfree,sec:multi}. They give an explicit bridge between the two supplied coefficient settings. The rational pairing in \cref{sec:algebraic} supplies an alternative when residue-level separation fails. These are proposed extensions of the supplied presentation, not a claim that their underlying coefficientwise or residue-theoretic ideas are absent from the literature.

\begin{center}\small
\begin{tabular}{@{}L{0.22\textwidth}L{0.35\textwidth}L{0.33\textwidth}@{}}
\toprule
Question & Successful formulation & Essential qualification\\
\midrule
Jordan separation & Definable circle in a real closed plane & Definably connected components\\[0.35em]
Coarse Jordan separation & Standard-part preimage of an ordinary Jordan curve & A thick wall in a non-Hausdorff topology\\[0.35em]
Cauchy vanishing & Closed Hahn-holomorphic $1$-form paired with an admissible null-homologous contour & Coherent coefficient data or a specified algebraic pairing\\[0.35em]
Generalized Stokes & Hahn forms on ordinary or infinitesimally deformed chains & Support-controlled pullback and integration\\[0.35em]
Radius-free residues & Microscopic torus after positive monomial rescaling & A determinant transform for a general complete intersection\\
\bottomrule
\end{tabular}
\end{center}

\section{Fields, summation, and three different topologies}\label{sec:foundations}

\subsection{Set-sized workspaces}
Let $\Gamma\subset\No$ be a nonzero divisible set-sized ordered additive subgroup. Write
\begin{equation}\label{eq:fields}
 R_\Gamma=\R((t^\Gamma)),\qquad K=\C((t^\Gamma))=R_\Gamma[i],\qquad
 t^\gamma=\omega^{-\gamma}.
\end{equation}
The Hahn-field theorem makes $K$ algebraically closed and $R_\Gamma$ real closed; the algebraic-closure statement is, for example, Corollary~4 of Poonen \cite{Poonen}. The normal-form interpretation places these fields inside $\SC$. The surreal background and restricted analytic evaluation are discussed by van den Dries and Ehrlich \cite{vdDE}. Every finite collection of points and every set-sized coefficient presentation used here can be accommodated by enlarging $\Gamma$ and taking its divisible hull.

For $a\ne0$, $v(a)$ is its least exponent. Put
\[
 K^\circ=\{a:v(a)\ge0\},\qquad \m=\{a:v(a)>0\},\qquad v(0)=+\infty.
\]
Every finite point has a unique form $a=c+\epsilon$, with $c\in\C$ and $\epsilon\in\m$. The map $\st(a)=c$ is reduction modulo $\m$. For ordinary $U\subset\C$, write $U\sh=\st^{-1}(U)\subset K^\circ$. The same notation can be used in the full class when no fixed workspace is intended.

The modulus $|x+iy|=\sqrt{x^2+y^2}$ is valued in $R_\Gamma$, not in $\R$. It satisfies the usual algebraic triangle inequality. Later we will also mention a real-valued \emph{valuation norm}; that is a different function.

\subsection{Strong summability is algebra, not a hidden limit}
A family of Hahn series is strongly summable when the union of its supports is well ordered and a fixed exponent receives contributions from only finitely many family members. We use the Hahn--Neumann support lemma \cite{Neumann,Poonen}: if $S,T$ are well ordered, then $S+T$ is well ordered with finite decompositions; if $E\subset\Gamma_{>0}$ is well ordered, its additive monoid $E^*$ is well ordered and a fixed exponent has only finitely many representations by finite words in $E$. Consequently the family indexed by
\[
 s+e_1+\cdots+e_q,\qquad s\in S,\quad e_j\in E,
\]
has the finiteness needed for coefficientwise regrouping. In particular, $(1+A)^{-1}=\sum_{q\ge0}(-A)^q$ for positive-support $A$.

A complex-linear map $L:V\to W$ extends to Hahn families by
\[
 L\left(\sum_\gamma v_\gamma t^\gamma\right)=\sum_\gamma L(v_\gamma)t^\gamma.
\]
It cannot enlarge support. No norm continuity is required. This observation will apply to ordinary integration, differentiation, residue extraction, and chosen linear splittings of differential complexes.

For $\Gamma=\Q+\Q\omega$, the strong sum $\sum_{n\ge0}t^n=(1-t)^{-1}$ exists, but its finite partial sums do not converge in the valuation topology: the tail valuations never exceed $\omega$. Strong summability therefore cannot be replaced by the statement that the valuation increases by a fixed positive amount at each step.

\subsection{The full fine topology and its restriction}
The full fine neighborhoods in $\SC$ are balls of all positive surreal radii. The following conclusions are proved in the analysis manuscript; we repeat the short argument because their exact scope is crucial.

\begin{proposition}[Full-class discreteness of sets]\label{prop:fullfine}
Every set-sized subset of $\SC$ is closed and discrete in the full fine topology. Every convergent net with a set-sized index set is eventually equal to its limit. In particular, the topology induced from the full fine topology on a fixed set-sized field $K$ is discrete.
\end{proposition}
\begin{proof}
For a point $p$ and a set $A$, the nonzero numbers $|p-a|$, $a\in A$, form a set of positive surreals. A surreal cut supplies a positive $\delta$ smaller than all of them. Thus $B(p,\delta)$ misses $A\setminus\{p\}$. This proves closedness and discreteness. Apply the same argument to the set of nonzero distances of the values of a net from its proposed limit. Eventual membership in the resulting ball forces eventual equality.
\end{proof}

This is not a theorem that a Hahn field with its \emph{own} valuation topology is discrete. The radius supplied by the cut may lie outside the workspace.

\subsection{The internal topology of a fixed Hahn field}
The internal ordered-field topology on $R_\Gamma^2=K$ uses radii from $R_\Gamma$. It agrees with the natural valuation topology. Indeed, monomial valuation balls refine ordered balls, and sufficiently small ordered monomial balls refine any prescribed valuation ball. For example, $v(x)>v(r)$ implies $|x|<r$ for $r>0$, while $|x|<t^{\gamma+\eta}$, $\eta>0$, implies $v(x)>\gamma$.

\begin{proposition}[Internal total separation]\label{prop:internal}
The internal topology on $K$ is nondiscrete and totally separated. Every ordinarily connected subset is a point. Every continuous map from an ordinary real interval or ordinary circle to $K$ is constant.
\end{proposition}
\begin{proof}
Every neighborhood of zero contains a nonzero monomial of sufficiently large exponent, so the topology is nondiscrete. Each valuation ball $\{z:v(z-a)>\gamma\}$ is clopen: a point outside it has a sufficiently small ball disjoint from it, by the ultrametric inequality. For $a\ne b$, the ball with $\gamma=v(b-a)$ contains $a$ and not $b$. Such a clopen separation splits any subset containing both points. Continuous images of connected ordinary spaces are connected.
\end{proof}

The same clopen separation idea works for the full fine topology, using $a+(b-a)\m$ with the full infinitesimal class. There is therefore no ordinary topological embedding $S^1(\R)\hookrightarrow\SC$ in that topology. In a fixed internally topologized $K$, every ordinary singular simplex is constant, so ordinary singular homology in positive degrees vanishes. The nonzero periods constructed below do not come from that singular chain complex.

\subsection{The residue topology}
On $K^\circ$ one may instead declare the open sets to be exactly $\st^{-1}(V)$ for ordinary open $V\subset\C$. This is the initial topology of $\st$. It has indiscrete fibers and is not Hausdorff, but its quotient by indistinguishability is exactly the ordinary complex plane. It deliberately forgets all infinitesimal distinctions. Its role is limited but clean: it supports a coarse Jordan theorem, not microscopic separation.

\section{Jordan separation: what survives and in which sense}\label{sec:jordan}

\subsection{A thick-wall Jordan theorem}
\begin{theorem}[Jordan separation in the residue topology]\label{thm:coarsejordan}
Let $J\subset\C$ be an ordinary Jordan curve, with bounded interior $D$ and exterior $E$. In the residue topology on $K^\circ$, the complement of
\[
 J\sh=\st^{-1}(J)
\]
has exactly two connected components, $D\sh$ and $E\sh$. Their common boundary is $J\sh$.
\end{theorem}
\begin{proof}
The ordinary Jordan theorem gives $\C\setminus J=D\sqcup E$, where both sets are nonempty and connected and both have boundary $J$. Every relatively open subset of $\st^{-1}(D)$ is the inverse image of an ordinary relatively open subset of $D$. A disconnection of $D\sh$ would consequently give a disconnection of $D$, since every fiber is nonempty and is wholly on one side of any such separation. Thus $D\sh$ is connected; the same proof applies to $E\sh$. They are disjoint relatively clopen subsets of the complement, so they are exactly its components. Finally, closure in this topology commutes with taking inverse images under the surjective standard-part map. Taking the boundaries proves the last assertion.
\end{proof}

The removed object is the entire monad of every point of $J$, not a one-dimensional embedded curve. A continuous coarse path whose endpoints have standard parts on different sides meets this thick wall. However, two infinitesimally separated points in the same boundary monad cannot be assigned different sides by this construction. Moreover, ``exterior'' here means that its standard-part image is ordinary unbounded: all of $K^\circ$ is bounded by any positive infinite element of $R_\Gamma$. We have not obtained an unbounded exterior in the internal ordered-field sense.

\subsection{The semialgebraic replacement}
A semialgebraic set over $R_\Gamma$ is described by finitely many polynomial equations and inequalities with coefficients in that field, and Boolean operations. A map is semialgebraic when its graph is. A set is \emph{semialgebraically connected} when it has no separation into two nonempty relatively open semialgebraic sets. Paths now have parameter domain
\[
 [0,1]_{R_\Gamma}=\{s\in R_\Gamma:0\le s\le1\},
\]
not just the ordinary real points in that interval. Definable connectedness and definable path connectedness agree for semialgebraic sets; this is part of the standard finite-topology theory \cite{BCR,BOhomology}.

\begin{theorem}[Established semialgebraic Jordan theorem]\label{thm:sajordan}
Let $j:S^1(R_\Gamma)\to R_\Gamma^2$ be a semialgebraic embedding, where
\[
 S^1(R_\Gamma)=\{(x,y):x^2+y^2=1\}.
\]
Then $R_\Gamma^2\setminus j(S^1(R_\Gamma))$ has exactly two semialgebraically connected components. One is bounded in $R_\Gamma^2$, the other is unbounded, and the common boundary of both is the image of $j$. Every semialgebraic continuous path between the two components meets that image.
\end{theorem}

This is an application of semialgebraic topology, not a new surreal theorem. One route uses compatible finite triangulation and the combinatorial separation theorem for an embedded polygonal circle on a triangulated sphere. The finite simplicial incidence data and connected-component calculations are independent of which real closed field realizes them. Alternatively, the definable Jordan--Brouwer theorem of Berarducci--Otero supplies the smooth-manifold case directly; the semialgebraic triangulation framework handles the topological formulation \cite{BCR,BOintersection,BOhomology}. The technical finite-topology input is being imported, not replaced by an informal appeal to transfer of arbitrary continuous maps.

Here is why finite triangulation has the right content. After putting the bounded image in a large box and compactifying the ambient plane semialgebraically, triangulate the pair consisting of the ambient sphere and the curve. The curve is a finite subcomplex which is a circle; the adjacent two-dimensional simplices split into two sides. Connectivity is computed by chains of incident simplices. The one component containing the point at infinity is the exterior. The other has bounded realization. Local incidence at the curve gives the common-boundary conclusion. A path missing the curve cannot move from one relatively open definable side to the other, because its parameter interval is definably connected.

The theorem does \emph{not} say that those components are connected in the unrestricted internal topology. By \cref{prop:internal}, all their ordinary connected components are singletons. The nondefinable cuts that disconnect the ordered field are excluded from the definition of semialgebraic connectedness.

\begin{example}[A microscopic circle with a genuine definable inside]
For $a\in K$ and $r\in R_\Gamma$, $r>0$, the circle
\[
 C_{a,r}=\{z\in K:|z-a|=r\}
\]
has semialgebraic inside $|z-a|<r$ and outside $|z-a|>r$. This remains true for $r=t^\rho$, for an infinite $r$, and for arbitrary $a$. Two points $a+r(1-t^\eta)$ and $a+r(1+t^\eta)$, $\eta>0$, lie on opposite definable sides although their normalized standard parts are equal.
\end{example}

\subsection{Winding and the limits of the analogy}
The semialgebraic circle has an integer-valued degree theory. The winding number of a loop around $a$ is the degree of its normalized direction $(\gamma-a)/|\gamma-a|$, or equivalently the corresponding definable homology invariant. With positive orientation, a simple circle has winding $1$ inside and $0$ outside. For piecewise polynomial loops Eisermann makes the construction effective through half of a Cauchy index, with a specified endpoint convention, computed by Sturm chains \cite[Theorem~1.2]{Eisermann}. We use the invariant rather than choosing a new sign convention for the index algorithm.

This degree is an ordinary integer, not a surreal number of turns. It provides a replacement for the part of complex contour theory that counts how a loop surrounds a point. It does not, by itself, define the integral of a general differential form. In particular, the primitive of a semialgebraic function need not be semialgebraic, and a chosen analytic expansion of a real closed field does not automatically contain all the radius-free Hahn germs of \cite{Geometry}.

A final distinction will recur: the set
\[
 \{a+re^{i\theta}:\theta\in[0,2\pi]_{\R}\}
\]
used as the trace of an ordinary-parameter contour is generally only a small subset of $C_{a,r}$. Its parameterization is not fine-continuous. The semialgebraic full circle, a coarse thick wall, and this contour trace are three different objects, even when they are denoted informally by the same equation.

\section{Why direct Riemann integration and unrestricted Stokes fail}\label{sec:obstructions}

\subsection{Finite meshes cannot resolve infinitesimal tolerances}
\begin{proposition}[Finite-partition obstruction]\label{prop:mesh}
Let $0=x_0<x_1<\cdots<x_N=1$ be a partition in $R_\Gamma$ with an ordinary finite number $N$ of subintervals. Its mesh is at least $1/N$. For $f(x)=x$, its right and left endpoint sums differ by
\begin{equation}\label{eq:riemannobstruction}
 \sum_{j=1}^N(x_j-x_{j-1})^2\ge\frac1N.
\end{equation}
Consequently ordinary finite tagged partitions cannot make all tag choices agree to a fixed positive infinitesimal error on $[0,1]$.
\end{proposition}
\begin{proof}
The increments $\Delta_j=x_j-x_{j-1}$ are positive and sum to one. At least one is at least $1/N$. The difference of the two sums is $\sum\Delta_j^2$. The algebraic identity
\[
 N\sum_j\Delta_j^2-\left(\sum_j\Delta_j\right)^2
 =\sum_{j<k}(\Delta_j-\Delta_k)^2\ge0
\]
proves \eqref{eq:riemannobstruction}. For any positive infinitesimal $\epsilon$, $1/N>2\epsilon$ for every ordinary $N$, so the two tag choices cannot both lie within $\epsilon$ of one proposed integral.
\end{proof}

There is a related logical pitfall. A definition that says ``all partitions of mesh less than $\delta$ have sums close to $I$'' becomes vacuous if $\delta$ is infinitesimal: there are no such ordinary finite partitions. One must specify an actual directed approximation system and prove it has enough elements. A polynomial antiderivative still gives $\int_0^1x\dd x=1/2$; the obstruction is to a particular unrestricted finite-mesh definition, not to all integration.

Hyperfinite partitions in a nonstandard universe can be useful, but they require additional structure, transfer, and a chosen relation between that universe and the surreal or Hahn workspace. They are not supplied merely by the field operations in $\No$. Taking only the ordinary standard part of such an integral would also discard the infinitesimal information sought here.

\subsection{The fundamental theorem already fails for arbitrary fine germs}
The class of infinitesimals and each of its additive cosets are fine-open and fine-closed. Therefore
\[
 q(z)=\begin{cases}0,&z\in\m,\\1,&z\notin\m\end{cases}
\]
is locally constant and hence fine-locally analytic, with $q'=0$. Yet $q(1)-q(0)=1$. No linear integration rule with integral zero for the zero integrand can satisfy the fundamental theorem
\[
 \int_\gamma dq=q(\gamma(b))-q(\gamma(a))
\]
for this $q$ and an admitted chain connecting $0$ to $1$. Thus a proposed Stokes theorem must restrict either forms, chains, or the differential under consideration. The Hahn theory restricts coefficient coherence, and the semialgebraic theory restricts definability.

The analysis manuscript gives an even sharper diagnostic: a globally fine-analytic logarithm $\mathcal L$ on $\SC^\times$ satisfying $\mathcal L'=1/z$ \cite[section ``A global fine-analytic logarithm'']{Analysis}. Nevertheless its Hahn unit-circle period is $2\pi i$. The restriction of $\mathcal L$ to ordinary points has the selected argument jump and is not a coherent primitive. Applying Stokes to $d\mathcal L$ on a coherent ordinary circle would mix two incompatible form categories.

\subsection{A contour avoiding every pole need not have an admissible pullback}
Consider the ordinary unit-circle trace and $a=1+t$ with $t>0$ infinitesimal. The function $1/(z-a)$ has no pole on that trace. Away from $z=1$ its Hahn expansion begins with $1/(z-1)$, but at $z=1$ its value is $-t^{-1}$. It cannot have a single Hahn expansion with continuous coefficient functions on the entire ordinary circle: the negative-exponent coefficient would vanish at all other ordinary points but not at $1$.

For $a=1-t$ the same problem occurs. Semialgebraic geometry distinguishes these poles as outside and inside the full unit circle, respectively. A coefficientwise contour integral based on the original macroscopic parameterization does not automatically do so. One needs a boundary-layer change of chart, a more general integration operator, or the algebraic pairing developed in \cref{sec:algebraic}. Avoidance of actual poles is necessary, but uniform admissibility of the pulled-back coefficients is an additional hypothesis.

\section{A Hahn de Rham complex and generalized Stokes}\label{sec:stokes}

\subsection{Forms and the chain pairing}
Let $M$ be an ordinary smooth manifold and let $\Omega^k(M;\C)$ denote its ordinary smooth complex-valued differential forms. Define
\begin{equation}\label{eq:hahnforms}
 \Omega_H^k(M)=\Omega^k(M;\C)((t^\Gamma)).
\end{equation}
An element is a family $\alpha=\sum_{\lambda\in S}\alpha_\lambda t^\lambda$ with a common well-ordered support $S$. This is a space of coefficient data over the ordinary manifold; it is not defined as the collection of all fine-smooth forms on a surcomplex point set. A derivative acts on $\alpha_\lambda$, not on $t^\lambda$.

Put
\[
 d_H\alpha=\sum_\lambda(d\alpha_\lambda)t^\lambda,\qquad
 \alpha\wedge\beta=\sum_\nu\left(\sum_{\lambda+\mu=\nu}
 \alpha_\lambda\wedge\beta_\mu\right)t^\nu.
\]
The support lemma makes the wedge product well defined. It immediately gives
\[
 d_H^2=0,\qquad
 d_H(\alpha\wedge\beta)=d_H\alpha\wedge\beta+(-1)^k\alpha\wedge d_H\beta
\]
for $\alpha$ of degree $k$. Ordinary smooth pullback is coefficientwise and commutes with $d_H$.

For a compact oriented ordinary $k$-manifold $C$ with corners and a smooth map $\phi:C\to M$, define
\begin{equation}\label{eq:chainintegral}
 \HInt_{(C,\phi)}\alpha
 =\sum_{\lambda\in S}t^\lambda\int_C\phi^*\alpha_\lambda.
\end{equation}
Its support is contained in $S$. Extend by finite linearity to piecewise smooth chains. Finite subdivisions, orientation reversal, ordinary reparameterization, and cancellation of oppositely oriented faces have their usual meanings. Noncompact chains could also be admitted with suitable coefficient support conditions, but compact chains suffice for all theorems stated here.

\begin{theorem}[Hahn-valued generalized Stokes]\label{thm:stokes}
Let $C$ be a compact oriented ordinary $k$-manifold with corners, with its boundary orientation given by the outward-normal-first convention. For $\alpha\in\Omega_H^{k-1}(M)$ and a smooth map $\phi:C\to M$,
\begin{equation}\label{eq:stokes}
 \boxed{\qquad
 \HInt_{(\partial C,\phi|_{\partial C})}\alpha
 =\HInt_{(C,\phi)}d_H\alpha.
 \qquad}
\end{equation}
The pairing is $K$-linear in the form and additive in the chain. The same identity holds for finite piecewise smooth chains.
\end{theorem}
\begin{proof}
For each exponent $\lambda$, ordinary Stokes gives
\[
 \int_{\partial C}(\phi|_{\partial C})^*\alpha_\lambda
 =\int_C d(\phi^*\alpha_\lambda)
 =\int_C\phi^*(d\alpha_\lambda).
\]
The two Hahn sums thus have equal coefficients at every exponent, which proves \eqref{eq:stokes}. Multiplication by $c=\sum c_\mu t^\mu\in K$ involves only finitely many pairs $\lambda+\mu=\nu$ at each target exponent. Ordinary complex-linearity of integration can therefore be applied to each finite convolution, proving $K$-linearity. For a finite chain apply the theorem to its pieces; internal boundary faces cancel.
\end{proof}

The ordinary Stokes theorem, including its chain and orientation conventions, is the classical input \cite{Lee}. The contribution of \cref{thm:stokes} is the precise Hahn coefficient setting and its compatibility with strong summability. No interchange of an analytic infinite limit with an integral occurs.

\subsection{Cohomology and periods}
\begin{theorem}[Cohomology comparison]\label{thm:cohomology}
There is a canonical isomorphism
\begin{equation}\label{eq:cohomology}
 H^k(\Omega_H^\bullet(M),d_H)
 \simeq H^k_{\mathrm{dR}}(M;\C)((t^\Gamma)).
\end{equation}
In particular, a closed Hahn form is exact if and only if each of its ordinary coefficients is exact. When ordinary de Rham cohomology is finite-dimensional, the right side is naturally
\[
 K\otimes_\C H^k_{\mathrm{dR}}(M;\C).
\]
\end{theorem}
\begin{proof}
Write $Z^k=\ker d$ and $B^k=\operatorname{im}d$ in the ordinary complex. A Hahn form is closed exactly when every coefficient belongs to $Z^k$. The ordinary quotient map $Z^k\to Z^k/B^k$ induces a map of closed Hahn families to $(Z^k/B^k)((t^\Gamma))$.

Choose a complex-linear section of that quotient. Its coefficientwise extension shows surjectivity without enlarging support. Choose also a complex-linear section $B^k\to\Omega^{k-1}(M;\C)$ of $d$. If all coefficient classes vanish, applying that section coefficientwise gives a Hahn primitive with the same containing support. The kernel is therefore precisely the exact Hahn forms. Although the proof uses linear sections, the induced isomorphism is the canonical coefficient-class map. In finite dimension, expansion in any finite basis identifies Hahn families with scalar extension by $K$.
\end{proof}

The finite-dimensional qualification matters. For an infinite-dimensional vector space, $V((t^\Gamma))$ generally contains families whose coefficients span infinitely many independent directions, whereas $K\otimes_\C V$ consists of finite sums of pure tensors. One must not silently replace the first object by the second.

On an ordinary star-shaped domain, the ordinary radial homotopy operator gives a primitive coefficientwise for every closed Hahn form of positive degree. On an arbitrary ordinary manifold, the classical de Rham theorem translates \eqref{eq:cohomology} into a period criterion: a closed Hahn form is exact exactly when all its periods over ordinary smooth cycles vanish \cite{Lee}. On $\C^\times$, for instance,
\[
 H^1(\Omega_H^\bullet(\C^\times))\simeq K,
\]
with generator represented by $dz/z$ and unit-circle period $2\pi i$. This nonzero cohomology is perfectly consistent with the vanishing of positive-dimensional \emph{ordinary singular} homology of the internally topologized point field $K$: the complexes have different bases and different chains.

\subsection{Positivity and an order-valued estimate}
Suppose $A(s)=\sum a_\lambda(s)t^\lambda$ has continuous real coefficients on an ordinary compact interval and $A(s)\ge0$ for every ordinary $s$. If it is not identically zero, its first nonzero coefficient function is nonnegative everywhere and positive somewhere; its ordinary integral is therefore positive. Thus
\[
 \HInt A(s)\dd s\ge0.
\]
For an ordinary piecewise $C^1$ path $\gamma$ of ordinary length $L$, and a uniformly supported continuous Hahn family $F$ along it, one obtains
\begin{equation}\label{eq:ml}
 \left|\HInt_\gamma F(z)\dd z\right|\le LM
 \quad\text{if }|F(\gamma(s))|\le M\in R_\Gamma,\ M\ge0.
\end{equation}
Indeed, for nonzero $I=\HInt_\gamma F\dd z$ use the fixed scalar $u=\overline I/|I|$ and apply positivity to
\[
 M|\gamma'(s)|-\operatorname{Re}\bigl(uF(\gamma(s))\gamma'(s)\bigr).
\]
Its integral is $LM-|I|$. This is the estimate proved in the analysis manuscript \cite{Analysis}; it also shows why a coefficientwise construction can retain an exact surreal, rather than merely standard-part, inequality.

\section{Infinitesimally deformed chains and homotopies}\label{sec:deformed}

\subsection{A support-controlled pullback}
Let $U\subset\R^d$ be ordinary open. Let $C$ be an ordinary compact smooth parameter manifold with corners, and suppose
\begin{equation}\label{eq:deformedmap}
 \Phi=\phi_0+\eta,\qquad
 \eta=\sum_{e\in E}\eta_e t^e,
\end{equation}
where $\phi_0:C\to U$ is smooth, $E\subset\Gamma_{>0}$ is well ordered, and the coordinate functions of $\eta_e$ are real-valued smooth functions on $C$. The value of $\eta$ at each ordinary parameter point is infinitesimal. The coefficient maps need not satisfy a uniform ordinary norm bound.

We use differential forms on $U$ whose Hahn coefficients are ordinary \emph{real-analytic} forms. For a coefficient function $a$, define its evaluation along $\Phi$ by its Taylor expansion:
\begin{equation}\label{eq:formalpullback}
 a(\phi_0+\eta)
 =\sum_{\alpha\in\N^d}\frac{(\partial^\alpha a)\circ\phi_0}{\alpha!}\eta^\alpha,
 \qquad d\Phi_j=d\phi_{0,j}+d\eta_j.
\end{equation}
For $\alpha=\sum_{\lambda,I} a_{\lambda,I}t^\lambda dx_I$, form $\Phi^*\alpha$ by substitution in the coefficients and the differentials. We use smooth forms on the \emph{parameter} manifold for the resulting coefficients. Analyticity is only required for the ambient coefficient functions whose infinitesimal evaluations are being interpreted.

\begin{lemma}[Admissible pullback]\label{lem:pullback}
If the ambient support is $S$, then \eqref{eq:formalpullback} defines
\[
 \Phi^*\alpha\in\Omega_H^\bullet(C),\qquad
 \supp(\Phi^*\alpha)\subset S+E^*.
\]
It respects wedge products, restrictions, and the differential:
\begin{equation}\label{eq:pullbackd}
 d_H\Phi^*\alpha=\Phi^*d_H\alpha.
\end{equation}
It also respects composition whenever the intervening ordinary coefficient maps are analytic and the same positive-support substitution hypotheses hold.
\end{lemma}
\begin{proof}
A contribution to one exponent is indexed by a source exponent in $S$ and a finite word of error exponents in $E$. Exterior products introduce only finitely many coordinate indices. The support lemma gives a well-ordered total support and only finitely many contributing words at every exponent. Hence each output coefficient is a finite sum of products of ordinary smooth functions and derivatives of ordinary analytic functions, and is smooth.

The ordinary product and chain rules prove \eqref{eq:pullbackd} formally. To justify their use, fix a target exponent, retain the finitely many coefficient functions and Taylor degrees that can contribute, and check the identity as an ordinary finite Taylor-jet identity. No discarded higher degree can affect that exponent. The same argument proves compatibility with products, restrictions, and analytic composition. Equality at every exponent proves the Hahn identities.
\end{proof}

The notation $\Phi$ describes an admissible parameterized chain, not a fine-continuous map from the ordinary parameter manifold. In particular, a nonconstant $\phi_0$ already prevents such continuity. This definition does not claim a canonical prolongation of every arbitrary smooth function on every fine domain.

\begin{definition}[Integral on a deformed chain]
For \eqref{eq:deformedmap} and an admissible form $\alpha$ of degree $\dim C$, put
\[
 \HInt_{(C,\Phi)}\alpha:=\HInt_C\Phi^*\alpha.
\]
Boundary parameterizations are obtained by restriction to the ordinary faces of $C$.
\end{definition}

\begin{theorem}[Stokes for Hahn-deformed chains]\label{thm:deformedstokes}
For a real-analytic-coefficient Hahn form $\alpha$ of degree $\dim C-1$,
\begin{equation}\label{eq:deformedstokes}
 \boxed{\qquad
 \HInt_{(\partial C,\Phi|_{\partial C})}\alpha
 =\HInt_{(C,\Phi)}d_H\alpha.
 \qquad}
\end{equation}
This integral is invariant under ordinary orientation-preserving smooth changes of parameter and finite admissible subdivision.
\end{theorem}
\begin{proof}
By \cref{lem:pullback}, the pulled-back form is a well-defined smooth Hahn form on $C$ and commutes with $d_H$. Apply \cref{thm:stokes} on $C$. Restriction to a face is the same as pulling back by the restricted parameterization. Ordinary change of variables at every coefficient proves reparameterization invariance; subdivision follows by finite boundary cancellation.
\end{proof}

One cannot replace ``ordinary changes of parameter'' by arbitrary fine-analytic reparameterizations without proving a separate substitution theorem and checking that the new pulled-back family remains admissible. The result does cover all the ordinary reparameterizations used in classical contour calculus and the explicitly controlled infinitesimal changes described next.

\subsection{The homotopy formula}
Let $\Psi:[0,1]_{\R}\times C\to U\sh$ be a Hahn-deformed map of the same kind, and write
\[
 \Psi^*\alpha=du\wedge\beta(u)+\gamma(u)
\]
with $\beta,\gamma$ forms in the $C$ variables. Define $\mathcal K\alpha=\HInt_0^1\beta(u)\dd u$.

\begin{theorem}[Hahn homotopy identity]\label{thm:homotopy}
For every admissible form,
\begin{equation}\label{eq:homotopy}
 \Psi_1^*\alpha-\Psi_0^*\alpha
 =d_H\mathcal K\alpha+\mathcal Kd_H\alpha.
\end{equation}
Consequently a closed form has equal integrals on two admissibly homotopic cycles.
\end{theorem}
\begin{proof}
The support bound of \cref{lem:pullback} is uniform over the ordinary compact parameter cylinder. The ordinary homotopy calculation, obtained by splitting off $du$ and using the fundamental theorem in $u$, holds at each Hahn exponent. Integrating a closed form over the difference of the two cycles then gives the integral of an exact form; Stokes on the parameter cycle makes it zero. Equivalently apply \cref{thm:deformedstokes} directly to the cylinder, with the usual opposite orientations on its two ends.
\end{proof}

For nonclosed forms there is a nonzero swept-region contribution. For example, a moving boundary applied to $\overline z\dd z$ changes the integral through $d(\overline z\dd z)=2i\dd x\wedge\dd y$. It would be incorrect to claim that all contour integrals are infinitesimally deformation invariant. Closedness, not merely existence of a derivative, is the controlling condition.

\section{Cauchy theory as a period theory}\label{sec:cauchy}

\subsection{Coherent holomorphy and the vanishing theorem}
Let $U\subset\C$ be an ordinary domain and
\[
 F(z)=\sum_{\lambda\in S}f_\lambda(z)t^\lambda\in\HH(U),
 \qquad f_\lambda\in\OO(U).
\]
At $z=c+\epsilon\in U\sh$, evaluation means
\begin{equation}\label{eq:coherentevaluation}
 F\sh(c+\epsilon)=\sum_{\lambda,n\ge0}
 t^\lambda\frac{f_\lambda^{(n)}(c)}{n!}\epsilon^n.
\end{equation}
The support is contained in $S+(\supp\epsilon)^*$, so this is strongly defined. It agrees with the supplied manuscript's notation and differential convention. In real coordinates,
\begin{equation}\label{eq:dFdz}
 d_H(F\dd z)=2i\,\overline\partial_HF\dd x\wedge\dd y=0.
\end{equation}
The last equality is coefficientwise ordinary holomorphy.

\begin{theorem}[Coherent Cauchy theorem]\label{thm:cauchy}
Let $F\in\HH(U)$. If $\gamma_0$ is an ordinary piecewise smooth closed cycle null-homologous in $U$, then
\[
 \HCoint_{\gamma_0}F\dd z=0.
\]
The same conclusion holds for a Hahn-deformed closed contour $\gamma=\gamma_0+\eta$ of the kind defined in \cref{sec:deformed}. More generally, its period equals the period of $\gamma_0$, whether or not that period is zero.
\end{theorem}
\begin{proof}
For the ordinary contour, apply ordinary Cauchy to every coefficient. Alternatively use a smooth filling chain and \eqref{eq:dFdz} in \cref{thm:stokes}. For the deformed contour, the homotopy $\gamma_u=\gamma_0+u\eta$, $0\le u\le1$, is admissible with the same positive error support. Apply \cref{thm:homotopy} to the closed form $F\dd z$.
\end{proof}

The null-homology condition can be replaced by simple connectedness of $U$ when discussing all closed contours. It cannot be omitted for a single closed contour in an arbitrary domain. Neither the contour nor a filling chain must be embedded, and the Jordan theorem is not needed for this vanishing statement. The Jordan theorem is useful when identifying a particular simple boundary with its interior, or when turning winding into an inside/outside count.

\subsection{An exact endpoint correction for a moving contour}
For $c\in U$ and $\epsilon\in\m$, define
\begin{equation}\label{eq:endpointT}
 T_F(c,\epsilon)=\sum_{n\ge0}
 \frac{D^nF(c)}{(n+1)!}\epsilon^{n+1}.
\end{equation}
It is the infinitesimal primitive increment, with no choice of global primitive.

\begin{theorem}[Exact contour-displacement formula]\label{thm:displacement}
Let $\gamma_0:[a,b]_{\R}\to U$ be piecewise smooth and let $\gamma=\gamma_0+\eta$ have positive-support smooth Hahn error, piecewise on the same finite partition. Then
\begin{align}\label{eq:displacement}
 \HInt_\gamma F\dd z
 &=\HInt_{\gamma_0}F\dd z
 +T_F(\gamma_0(b),\eta(b))
 -T_F(\gamma_0(a),\eta(a)).
\end{align}
In particular, a closed deformation has exactly the original period, not just the same standard part.
\end{theorem}
\begin{proof}
Set $\gamma_u(s)=\gamma_0(s)+u\eta(s)$. Formal differentiation gives
\[
 \partial_u\bigl(F(\gamma_u)\partial_s\gamma_u\bigr)
 =\partial_s\bigl(F(\gamma_u)\partial_u\gamma_u\bigr).
\]
All products and derivatives are admissible by \cref{lem:pullback}. Integrate coefficientwise over the ordinary rectangle $[0,1]\times[a,b]$. The right side gives the two endpoint terms
\[
 \HInt_0^1 F(\gamma_0(b)+u\eta(b))\eta(b)\dd u
 -\HInt_0^1 F(\gamma_0(a)+u\eta(a))\eta(a)\dd u.
\]
Taylor expansion and the ordinary integral of $u^n$ identify these as \eqref{eq:endpointT}. Splitting into finitely many smooth pieces introduces endpoint terms which cancel at the joins. For a closed contour the remaining two endpoint terms coincide.
\end{proof}

This theorem extends the endpoint construction in the analysis manuscript: it shows that every admissible infinitesimal deformation of the entire path gives exactly that construction. It holds on nonsimply connected $U$. A primitive is not required because only the homotopy strip, not an arbitrary filling disk, is used.

\subsection{Winding and the Cauchy kernel at a displaced point}
Let $a=a_0+\delta$ with $a_0\in\C$, $\delta\in\m$, and suppose $a_0$ is not on the ordinary contour $\gamma_0$. On a common ordinary collar of the contour,
\[
 \frac1{z-a}=\sum_{n\ge0}\frac{\delta^n}{(z-a_0)^{n+1}}.
\]
This is an admissible Hahn family. Higher-pole terms have zero periods, and \cref{thm:displacement} permits an admissible deformation of the contour. Therefore
\begin{equation}\label{eq:wind}
 \frac1{2\pi i}\HCoint_{\gamma_0+\eta}\frac{\dd z}{z-a}
 =\Wind(\gamma_0,a_0)\in\Z.
\end{equation}
The hypothesis that $a_0$ misses the contour is precisely the one that fails for $a=1\pm t$ in \cref{sec:obstructions}.

\begin{theorem}[Cauchy formula with infinitesimal arguments]\label{thm:cauchyformula}
Let $D$ be a bounded ordinary Jordan domain with piecewise smooth positively oriented boundary and $\overline D\subset U$. For $z=c+\epsilon\in D\sh$ and $n\ge0$,
\begin{equation}\label{eq:cauchyformula}
 D^nF\sh(z)=\frac{n!}{2\pi i}
 \HCoint_{\partial D}\frac{F(\zeta)}{(\zeta-z)^{n+1}}\dd\zeta.
\end{equation}
The same value results from an admissibly infinitesimally deformed boundary.
\end{theorem}
\begin{proof}
Expand
\[
 (\zeta-c-\epsilon)^{-n-1}
 =\sum_{k\ge0}\binom{n+k}{k}
 \epsilon^k(\zeta-c)^{-n-k-1}.
\]
The ordinary point $c$ has positive ordinary distance from the compact boundary, so all coefficient functions are holomorphic on a common collar. The support lemma permits multiplication by $F$ and coefficientwise integration. Ordinary Cauchy gives the coefficient
$f_\lambda^{(n+k)}(c)/k!$, yielding \eqref{eq:coherentevaluation} for $D^nF$. The kernel form is closed on the collar, so the last assertion follows from \cref{thm:displacement}.
\end{proof}

On a simply connected ordinary domain, normalized ordinary primitives of all $f_\lambda$ assemble into a coherent primitive with support in $S$. On a general domain such a primitive exists exactly when all the periods vanish. Similarly, Morera's criterion holds for common-support continuous coefficient families: vanishing Hahn integrals around every ordinary triangle means vanishing of every ordinary coefficient integral, hence holomorphy by ordinary Morera. These recover the corresponding statements of \cite{Analysis} and locate them in the chain-and-form framework.

\subsection{Moving poles and actual zero clusters}
For a coherent fraction $F/G$, assume that the ordinary leading coefficient of $G$ has no zero on the contour and has finitely many zeros in the enclosed ordinary region. The preparation and moving-pole theorem of \cite{Analysis} then gives
\begin{equation}\label{eq:movingres}
 \frac1{2\pi i}\HCoint_{\partial D}\frac FG\dd z
 =\sum_{a\in D\sh}\Res_a\left(\frac FG\dd z\right),
\end{equation}
where the sum is over actual surcomplex poles, not merely their ordinary reductions. The sum is finite; its summands may be infinite. The same collar homotopy preserves the left side under an admissible boundary deformation. Applied to $DG/G$, it counts zeros minus poles with their ordinary integer multiplicities.

The algebraic mechanism behind \eqref{eq:movingres} is worth emphasizing. Locally prepare $G$ as a coherent unit times a polynomial with infinitesimal roots. Divide the numerator by that polynomial. The analytic quotient integrates to zero, and partial fractions identify the remaining contour integral with the actual local residues. This is stronger than calling a coefficientwise residue at the reduced center ``the residue'' and leaving its relation to the split points unspecified.

\section{Affine scales, microscopic contours, and compatible charts}\label{sec:scales}

\subsection{A contour at an arbitrary surreal scale}
Suppose a function $f$ on $a+rU\sh$, with $a\in K$ and $r\in R_\Gamma$, $r>0$, has a coherent chart
\[
 f(a+rw)=F\sh(w),\qquad F\in\HH(U).
\]
For an ordinary path $\gamma$ in $U$, define
\begin{equation}\label{eq:affineintegral}
 \HInt_{a+r\gamma}f(z)\dd z
 :=r\HInt_\gamma F(w)\dd w.
\end{equation}
This is the affine contour construction of \cite{Analysis}. It also works for a Hahn-deformed parameter path, provided the normalized pullback is admissible. In several real coordinates one inserts the ordinary algebraic Jacobian and wedge factors of the affine change.

Cauchy and Stokes transfer through this definition with their correct differential factors. For instance, a pole at $a+r\xi$ contributes the same residue after substitution, since the factor $r$ in $dz$ cancels the scaling of a simple-pole denominator. A $k$-dimensional uniform real scaling contributes $r^k$ to a $k$-form. Forgetting that factor would incorrectly erase microscopic areas and volumes.

\subsection{A precise compatibility criterion}
\begin{proposition}[Agreement on admissible overlaps]\label{prop:overlap}
Suppose two chart descriptions of the same parameterized oriented chain produce pointwise equal pulled-back differential forms on every ordinary parameter point, and both pulled-back forms are common-support Hahn families with smooth coefficients. Then their integrals agree. The same holds for a finite collection of chart pieces with matching oriented faces.
\end{proposition}
\begin{proof}
Enlarge the workspace to contain both supports. Uniqueness of Hahn normal forms makes each coefficient of the difference vanish at every parameter point. Thus the coefficient forms are equal. Their ordinary integrals agree, and so do their Hahn sums. For finitely many pieces use additivity and cancellation of matching faces.
\end{proof}

This criterion is intentionally stronger than equality of the set-theoretic images of two curves. A chain includes multiplicity, orientation, and the pulled-back differential. It is also weaker than requiring one global chart for every scale: finitely many verified overlap identities suffice. What is not justified is a patchwork of arbitrary local fine germs for which common-support pullbacks and their identities have not been established.

\subsection{Macroscopic and microscopic winding resolve different information}
A macroscopic contour whose reduction avoids a cluster assigns the same winding to every point in that cluster. A contour centered on one actual pole and rescaled below the finite separation distances can resolve an individual contribution. This does not require a surreal-valued winding number: ordinary integer winding numbers in different charts suffice.

For example, the poles of $e^z/(z^2-t)$ are $\pm t^{1/2}$. The ordinary unit circle surrounds both, and so does the microscopic circle of radius $t^{1/4}$ about zero. A circle centered at $t^{1/2}$ of radius $t$ isolates the positive pole. In the coordinate $z=t^{1/2}+tw$, the second pole is at $w=-2t^{-1/2}$, an infinite normalized distance. Each integral is calculated in its own admissible chart. The outer integral is the sum of the two inner residues, even though the individual residues are infinite. There is no assertion that these ordinary-parameter traces constitute ordinary topological circles in the fine plane.

\section{Radius-free germs become coherent after microscopic rescaling}\label{sec:radiusfree}

\subsection{The missing common-radius problem}
The geometry manuscript admits
\begin{equation}\label{eq:radiusfree-ring}
 \A_n=\C\{z_1,\ldots,z_n\}((t^\Gamma)).
\end{equation}
An element $H=\sum_\lambda h_\lambda(z)t^\lambda$ has analytic germs as coefficients. There is no requirement that all $h_\lambda$ extend to the same ordinary polydisk. For instance,
\begin{equation}\label{eq:shrinking}
 H(z)=\sum_{m\ge1}\frac{t^m}{1-mz}
\end{equation}
is legitimate, but its coefficient poles $1/m$ prevent a common ordinary disk of positive radius. An ordinary fixed-circle coefficient integral at zero would therefore be unavailable. Evaluation at infinitesimal points is nevertheless strongly defined.

\begin{theorem}[Positive rescaling gives polynomial coefficient families]\label{thm:polynomialization}
Let $\rho_1,\ldots,\rho_n\in\Gamma_{>0}$, put $r_j=t^{\rho_j}$, and write $rw=(r_1w_1,\ldots,r_nw_n)$. The Taylor substitution
\begin{equation}\label{eq:rescaling}
 \sigma_r(H)(w):=H(rw)
 =\sum_{\lambda,\alpha}h_{\lambda,\alpha}
 t^{\lambda+\alpha\cdot\rho}w^\alpha
\end{equation}
defines an injective $K$-algebra homomorphism
\[
 \sigma_r:\A_n\longrightarrow\C[w_1,\ldots,w_n]((t^\Gamma)).
\]
Every fixed Hahn coefficient in \eqref{eq:rescaling} is an ordinary polynomial. Its support is contained in $\supp H+\{\rho_1,\ldots,\rho_n\}^*$. It has coherent evaluation on the thickening of every ordinary domain, and
\[
 \partial_{w_j}\sigma_r(H)=r_j\sigma_r(\partial_{z_j}H).
\]
\end{theorem}
\begin{proof}
Write the Taylor germ $h_\lambda(z)=\sum_\alpha h_{\lambda,\alpha}z^\alpha$. At a fixed Hahn exponent $\nu$, possible contributions satisfy
\[
 \lambda+\alpha_1\rho_1+\cdots+\alpha_n\rho_n=\nu.
\]
The support lemma gives finitely many such source exponents and finite words in the positive alphabet $\{\rho_1,\ldots,\rho_n\}$. There are consequently finitely many multi-indices $\alpha$ contributing at $\nu$. The coefficient is a polynomial, and the total support is well ordered.

Taylor multiplication and differentiation are valid at each exponent because the relevant regroupings are finite. This gives the algebra and derivative statements. The Taylor embedding $\A_n\hookrightarrow K[[z]]$ is injective, and diagonal substitution $z_j=r_jw_j$ is injective in that formal ring since $r_j\ne0$. It follows that \eqref{eq:rescaling} is injective. Finally, polynomial coefficients are holomorphic on every ordinary domain. Evaluation at a finite normalized point agrees with evaluation of $H$ at the infinitesimal point $rw$, by the same finite-contribution Taylor identities.
\end{proof}

No coefficient radii enter this proof. In fact the argument works even for arbitrary formal complex coefficient series. Its analytically relevant conclusion is not ordinary convergence of the original germs on a uniform disk, but a \emph{new coherent representation in an infinitesimal chart}. It neither establishes nor requires all-scale coherence of the original function.

For \eqref{eq:shrinking}, substitution $z=t^{1/4}w$ gives
\[
 H(t^{1/4}w)=\sum_{m\ge1}\sum_{q\ge0}m^q t^{m+q/4}w^q.
\]
At a fixed exponent only finitely many $m,q$ occur. The singularities of the original ordinary coefficient representatives no longer obstruct coefficientwise contour integration in the $w$-plane.

\subsection{One-variable perturbation residues}
Let $f\in\C\{z\}$ have order $m\ge1$ at zero, and let
\[
 F=f+E\in\A_1,\qquad \vH(E)>0.
\]
Zero perturbations are permitted. For $H\in\A_1$ define, as in the geometry manuscript,
\begin{equation}\label{eq:residueone}
 \Res_F(H)=\sum_{k\ge0}(-1)^k\res_{f^{k+1}}(HE^k),
 \qquad
 \res_{f^{k+1}}(q)=\Res_{z=0}\frac{q(z)\dd z}{f(z)^{k+1}}.
\end{equation}
The ordinary residue operator acts coefficientwise on the Hahn numerator. The positive support of $E$ proves strong summability with output support in $\supp H+(\supp E)^*$.

\begin{theorem}[Microscopic circle realization]\label{thm:microcircle}
Assume $E\ne0$ and choose $\rho>0$ such that $m\rho<\vH(E)$. If $E=0$, choose any $\rho>0$. Put $r=t^\rho$. Then the microscopic circle functional is well defined and
\begin{equation}\label{eq:microcircle}
 \boxed{\qquad
 \frac1{2\pi i}\HCoint_{rS^1}\frac{H(z)}{F(z)}\dd z
 =\Res_F(H).
 \qquad}
\end{equation}
The integral means substitution $z=rw$ followed by coefficientwise integration over the ordinary unit circle in $w$. It is independent of every choice of $\rho$ satisfying the displayed bound.
\end{theorem}
\begin{proof}
Write $f(z)=cz^m+O(z^{m+1})$, $c\ne0$. By \cref{thm:polynomialization}, $H(rw)$ and $F(rw)$ have polynomial Hahn coefficients. Moreover,
\[
 r^{-m}F(rw)=cw^m+P(w),\qquad \vH(P)>0.
\]
Higher terms of $f$ acquire positive relative exponents, and the perturbation terms have relative exponent at least $\vH(E)-m\rho>0$. On an ordinary annulus around $|w|=1$, the leading coefficient $cw^m$ is nowhere zero. Geometric inversion therefore gives a common-support holomorphic Hahn family for $rH(rw)/F(rw)$ on that annulus.

Expand instead around $f(rw)$:
\[
 \frac{rH(rw)}{F(rw)}
 =\sum_{k\ge0}(-1)^k
 \frac{rH(rw)E(rw)^k}{f(rw)^{k+1}}.
\]
This is strongly summable on the annulus, since $E(rw)/f(rw)$ has strictly positive Hahn valuation there. Coefficientwise integration extracts the coefficient of $w^{-1}$. For an ordinary germ $q$ the formal Laurent change of variable gives
\[
 [w^{-1}]\frac{rq(rw)}{f(rw)^{k+1}}
 =[z^{-1}]\frac{q(z)}{f(z)^{k+1}}.
\]
The identity holds as well for a Hahn family of germs, by finite contributions after rescaling. Apply it to $q=HE^k$ and sum the strongly summable family. This gives exactly \eqref{eq:residueone}. Since the right side does not involve $\rho$, the last assertion follows.
\end{proof}

The bound is a \emph{protecting-radius} condition: the chosen circle must be infinitesimally large relative to the perturbation scales after accounting for the order $m$. For an infinitesimal zero $a$ of $F$, $v(f(a))=m v(a)$ while $v(E(a))\ge\vH(E)$; hence $v(a)\ge\vH(E)/m>\rho$. Thus every zero in the originating infinitesimal cluster lies strictly inside the protecting circle. The finite-zero and actual-residue interpretation follows from the one-variable preparation theorem, either before or after rescaling \cite{Analysis,Geometry}.

Radius independence here is an exact residue identity. It does not require a convergent sequence of radii or a fine-continuous homotopy between scales whose ratio might be infinite.

\section{Microscopic tori and general complete-intersection residues}\label{sec:multi}

\subsection{Coordinate-power leading systems}
Let
\begin{equation}\label{eq:powerleading}
 F_j(z)=z_j^{d_j}+E_j(z),\qquad d_j\ge1,\quad E_j\in\A_n,
 \quad \vH(E_j)>0
\end{equation}
for each nonzero $E_j$. Choose positive $\rho_j$ such that
\begin{equation}\label{eq:protecting}
 d_j\rho_j<\vH(E_j)\quad\text{whenever }E_j\ne0,
 \qquad r_j=t^{\rho_j}.
\end{equation}
These choices exist by divisibility of $\Gamma$. For a zero perturbation there is no upper restriction on $\rho_j$.

Write $T_r$ for the parameterized real $n$-torus
\[
 z_j=r_je^{i\theta_j},\qquad 0\le\theta_j\le2\pi,
\]
with ordinary parameter angles and orientation $d\theta_1\wedge\cdots\wedge d\theta_n$. Integration is defined by its pullback to the ordinary product of unit circles. It is a real $n$-dimensional cycle in an $n$-complex-variable calculation, not the $(2n-1)$-dimensional sphere used in some other several-variable integral formulas.

\begin{theorem}[Coordinate-power torus realization]\label{thm:microtorus}
For $H\in\A_n$, the following integral is well defined:
\begin{equation}\label{eq:microtorus}
 \mathcal R_F(H):=\frac1{(2\pi i)^n}
 \HInt_{T_r}\frac{H(z)\dd z_1\wedge\cdots\wedge\dd z_n}
 {F_1(z)\cdots F_n(z)}.
\end{equation}
It satisfies the exact coefficient formula
\begin{equation}\label{eq:toruscoeff}
 \boxed{\quad
 \mathcal R_F(H)=
 \sum_{k\in\N^n}(-1)^{|k|}
 [z_1^{d_1(k_1+1)-1}\cdots z_n^{d_n(k_n+1)-1}]
 \left(H\prod_{j=1}^nE_j^{k_j}\right).
 \quad}
\end{equation}
This is the perturbation residue of \cite{Geometry}. Its support is contained in $\supp H+E^*$, where $E=\bigcup_j\supp E_j$. It is independent of all protecting radii satisfying \eqref{eq:protecting}.
\end{theorem}
\begin{proof}
After substituting $z_j=r_jw_j$ we have
\[
 r_j^{-d_j}F_j(rw)=w_j^{d_j}+P_j(w),\qquad \vH(P_j)>0.
\]
All coefficient functions are polynomials by \cref{thm:polynomialization}. On a common ordinary product annulus about $|w_j|=1$, expand each reciprocal geometrically. The shifts of the error supports are positive by \eqref{eq:protecting}; the support lemma therefore proves strong summability of the entire pulled-back form.

At multi-index $k$ the denominator contributes $\prod_j z_j^{-d_j(k_j+1)}$, and the numerator is $H\prod_jE_j^{k_j}$. The ordinary torus integral selects the Laurent monomial $w_1^{-1}\cdots w_n^{-1}$. A source Taylor monomial $z^\alpha$ contributes only when
\[
 \alpha_j=d_j(k_j+1)-1\quad\text{for every }j.
\]
For that monomial the scale factor in coordinate $j$ is
\[
 r_j^{\,1+\alpha_j-d_j(k_j+1)}=1.
\]
The surviving coefficients are exactly \eqref{eq:toruscoeff}, including its sign. The support statement follows because the ordinary Taylor coefficient operators do not enlarge Hahn support and the perturbations have positive support. The formula contains no $r_j$, proving radius independence.
\end{proof}

The theorem permits different valuation ranks in different coordinates. For example, $\rho_1\in\Q_{>0}$ and $\rho_2$ proportional to $\omega$ are legitimate. Nothing in the proof replaces $\Gamma$ by a real-valued norm. It also supplies a contour meaning to a residue that was deliberately defined algebraically in the supplied geometry manuscript.

\subsection{Arbitrary isolated analytic complete intersections}
Suppose $f=(f_1,\ldots,f_n)$ is an ordinary analytic germ at zero with an isolated common zero:
\[
 0<\dim_\C\C\{z\}/(f_1,\ldots,f_n)<\infty.
\]
Let $F=f+E$, with all nonzero $E_j$ of strictly positive Hahn valuation. The perturbation residue from \cite{Geometry} is
\begin{equation}\label{eq:generalres}
 \Res_F(H)=\sum_{k\in\N^n}(-1)^{|k|}
 \res_{(f_1^{k_1+1},\ldots,f_n^{k_n+1})}
 \left(H\prod_jE_j^{k_j}\right).
\end{equation}
Here the ordinary local Grothendieck residue is a complex-linear functional on analytic germs. Its coordinate-power case is Taylor coefficient extraction. We use the classical residue transformation law, as presented in \cite{Tsikh,CDS}.

Since $(f)$ has finite colength, choose $d_1,\ldots,d_n\ge1$ and an ordinary analytic matrix $A(z)$ such that
\begin{equation}\label{eq:transformmatrix}
 A(z)f(z)=(z_1^{d_1},\ldots,z_n^{d_n})^{\mathsf T}.
\end{equation}
Set $G=AF$. Then
\[
 G_j=z_j^{d_j}+D_j,\qquad D=AE,
\]
and all nonzero $D_j$ have positive Hahn valuation. No invertibility of $A$ is assumed.

\begin{lemma}[Transformation of the Hahn perturbation residue]\label{lem:residuetransform}
With these hypotheses,
\begin{equation}\label{eq:residuetransform}
 \Res_F(H)=\Res_G(H\det A).
\end{equation}
\end{lemma}
\begin{proof}
We first specify the classical input. For a finite analytic parameter deformation $f_u$ of an isolated complete intersection, the sum of local residues over the nearby cluster is analytic in $u$. If $g_u=A f_u$ also has an isolated central fiber, the classical transformation law gives
\[
 \res^{\mathrm{cluster}}_{f_u}(h)
 =\res^{\mathrm{cluster}}_{g_u}(h\det A).
\]
The cluster for $g_u$ can contain additional points; the determinant is part of the law. Taylor expansion in $u$ gives the geometric denominator expansions appearing in \eqref{eq:generalres}. This ordinary finite-family statement follows from the usual local residue transformation law and its parameter version \cite{Tsikh,CDS}; the same finite-family use of residue theory is explained in the geometry manuscript.

Now fix a Hahn exponent $\nu$. In either side of \eqref{eq:residuetransform}, a coefficient is built from a source coefficient of $H$ and a finite word of positive perturbation exponents. The support lemma leaves only finitely many contributing words and hence finitely many ordinary coefficient germs. Replace the finitely many retained perturbation terms by independent ordinary parameters. Together with the finitely many entries of $A$, their germs have a common ordinary representative neighborhood. Apply the finite-family identity there, compare the finitely many required Taylor coefficients, and substitute the associated monomials $t^e$ for the parameters. Exponent collisions cause only finite additions. Thus the two sides have equal coefficient at $\nu$. Since $\nu$ was arbitrary, \eqref{eq:residuetransform} follows.
\end{proof}

\begin{theorem}[Microscopic contour representation of a general residue]\label{thm:generalrepresentation}
Choose $A,d_j$ as in \eqref{eq:transformmatrix} and protecting radii for $G$ as in \eqref{eq:protecting}. Then
\begin{equation}\label{eq:generalrepresentation}
 \boxed{\qquad
 \Res_F(H)=\frac1{(2\pi i)^n}
 \HInt_{T_r}
 \frac{H(z)\det A(z)\dd z_1\wedge\cdots\wedge\dd z_n}
 {G_1(z)\cdots G_n(z)}.
 \qquad}
\end{equation}
The value is independent of the admissible radii and of the choice of ordinary transformation data $A,d_j$.
\end{theorem}
\begin{proof}
The integrand on the right is admissible by \cref{thm:polynomialization,thm:microtorus}, which identify its integral with $\Res_G(H\det A)$. Apply \cref{lem:residuetransform}. Every permitted choice gives the same independently defined left side, proving independence.
\end{proof}

This is a contour \emph{representation by a transformed form}. It does not assert that the untransformed form $H\,dz_1\wedge\cdots\wedge dz_n/(F_1\cdots F_n)$ has the same integral over that coordinate torus. In the classical theory its natural Grothendieck cycle is adapted to the equations, typically $|F_j|=\varepsilon_j$. Replacing that cycle by a coordinate torus without the determinant transformation would in general be wrong. Formula \eqref{eq:generalrepresentation} is useful precisely because it gives a controlled coordinate-torus substitute.

Nor does the theorem claim that $F$ and $G$ have identical zero schemes. Multiplication by a matrix with nonunit determinant can introduce additional zeros and multiplicities. The determinant compensates for this in the residue identity; it is not an optional Jacobian factor.

\subsection{Relation to the finite-algebra conclusions of the supplied source}
The geometry manuscript proves that the perturbed quotient has the same finite dimension as the ordinary complete-intersection algebra, and that
\begin{equation}\label{eq:traceresidue}
 \Res_F(HJ_F)=\Tr(M_H)=\sum_a e_aH(a),\qquad
 J_F=\det\left(\frac{\partial F_i}{\partial z_j}\right).
\end{equation}
For simple zeros it gives $\Res_F(H)=\sum_a H(a)/J_F(a)$. These are reported companion results, not reproved here as part of the contour construction. Combining them with \cref{thm:generalrepresentation} identifies the same functional simultaneously as a finite-algebra trace pairing and as a microscopic integral. The contour realization itself only uses the strong residue expansion, ordinary finite-family residue transformation, and positive rescaling; it does not require the full Nullstellensatz or faithful-flatness theorems stated in that manuscript.

In dimension one, Stokes, the residue of a Laurent series, and the integer winding of a circle already organize the calculation. In higher dimensions the determinant transformation and the real $n$-torus play the corresponding roles. There is no need to postulate a new Jordan separation theorem for a $(2n-1)$-sphere in order to justify these particular residues.

\section{An algebraic contour pairing for rational differentials}\label{sec:algebraic}

\subsection{Definition from definable winding}
Let $\gamma:[0,1]_{R_\Gamma}\to K$ be a semialgebraic continuous loop, piecewise smooth when differential notation is used. For $a\notin\gamma$, write $\Wind_{\rm sa}(\gamma,a)\in\Z$ for its semialgebraic winding number. The normalization is $1$ for a positively oriented circle about an inside point and $0$ outside. These are the established degree invariants discussed in \cref{sec:jordan}.

For a rational differential $\omega=Q(z)\dd z$, $Q\in K(z)$, with no pole on $\gamma$, define
\begin{equation}\label{eq:algebraicpairing}
 \boxed{\qquad
 I_{\rm alg}(\gamma,\omega)
 :=2\pi i\sum_{a\in K}\Wind_{\rm sa}(\gamma,a)\Res_a\omega.
 \qquad}
\end{equation}
Only the finitely many finite poles contribute. The factor $2\pi i$ is the normalization selected to match ordinary complex integration; it is not forced by ordered-field arithmetic alone. A residue at infinity could be included after compactifying the curve and changing charts, but is unnecessary for this affine formula.

Equation \eqref{eq:algebraicpairing} is a \emph{definition of a pairing}, not a residue theorem claimed for a previously constructed Riemann integral. It is useful because it requires neither ordinary leading separation nor a global logarithm. Poles and contour coefficients may lie at arbitrary Hahn scales.

\begin{theorem}[Characterization and algebraic Cauchy theory]\label{thm:algpairing}
For a fixed semialgebraic loop, \eqref{eq:algebraicpairing} is the unique $K$-linear pairing on rational differentials without poles on the loop which vanishes on exact rational differentials and satisfies
\[
 I_{\rm alg}\left(\gamma,\frac{\dd z}{z-a}\right)
 =2\pi i\Wind_{\rm sa}(\gamma,a).
\]
It is additive under concatenation of loops and invariant under definable homotopies avoiding the poles. If $\gamma$ is a positively oriented semialgebraic Jordan curve and $Q$ has no poles in its bounded side or on its boundary, then
\[
 I_{\rm alg}(\gamma,Q\dd z)=0.
\]
If $a$ is on that bounded side, then
\[
 I_{\rm alg}\left(\gamma,\frac{Q(z)}{z-a}\dd z\right)=2\pi i Q(a).
\]
\end{theorem}
\begin{proof}
Since $K$ is algebraically closed of characteristic zero, partial fractions give
\[
 Q(z)=P(z)+\sum_a\sum_{j=1}^{m_a}\frac{c_{a,j}}{(z-a)^j}.
\]
The polynomial part has a polynomial primitive. For $j\ge2$,
\[
 \frac{\dd z}{(z-a)^j}
 =d\left(\frac{(z-a)^{1-j}}{1-j}\right).
\]
Thus a pairing with the required properties must retain exactly the $j=1$ terms, with the displayed winding coefficients. This proves uniqueness and also verifies that \eqref{eq:algebraicpairing} has those properties. Additivity and homotopy invariance follow from the corresponding winding properties. For a positive Jordan curve, exterior poles have winding zero. In the Cauchy-kernel expression the only contributing interior pole is $a$, with residue $Q(a)$.
\end{proof}

For a rational function $Q$ with neither zeros nor poles on the loop, logarithmic differentiation gives
\[
 \frac1{2\pi i}I_{\rm alg}\left(\gamma,\frac{Q'}Q\dd z\right)
 =\sum_a\ord_a(Q)\Wind_{\rm sa}(\gamma,a).
\]
This is a rational argument principle over the full real closed workspace. Eisermann's polynomial Cauchy and root-counting theorems provide the effective winding component for polynomially described contours \cite{Eisermann}. An arbitrary Hahn coefficient still needs an effective arithmetic and order representation before this becomes an implemented algorithm.

\subsection{The local algebra behind the normalization}
For a formal Laurent differential over $K$,
\[
 \omega=\sum_{j\ge j_0}c_j u^j\dd u\in K((u))\dd u,
\]
there is a formal Laurent primitive exactly when $c_{-1}=0$. Every other term integrates by division by $j+1$, whereas differentiating a Laurent series can never produce a $u^{-1}$ coefficient. Hence
\begin{equation}\label{eq:formalresclass}
 K((u))\dd u\,/\,dK((u))\simeq K,\qquad[\omega]\longmapsto c_{-1}.
\end{equation}
The residue is invariant under an invertible formal local coordinate change $u=v h(v)$, $h(0)\ne0$. It suffices to check exact differentials, which remain exact, and $du/u=dv/v+h'(v)\,dv/h(v)$, whose second term has zero residue. Thus residues give a coordinate-invariant obstruction to exactness independently of any connected topology.

Globally, partial fractions give the analogous statement for rational differentials: modulo exact rational differentials, their classes are the finite residue vectors. Formula \eqref{eq:algebraicpairing} pairs those vectors with integer winding vectors. It is an explicit algebraic model of the period pairing rather than a new topological connectedness assertion.

\subsection{Agreement with Hahn integration on a verified overlap}
\begin{proposition}[Comparison on separated charts]\label{prop:comparison}
Suppose a semialgebraic loop has an affine normalized description whose restriction to ordinary parameters is an admissible Hahn-deformed contour $\gamma_0+\eta$. Assume the same description on the full field interval is uniformly infinitesimally close to its ordinary real-coefficient leading loop. For each pole, suppose its normalized position is finite with reduction off $\gamma_0$, or is infinite. Whenever the rational pulled-back differential is a coherent family on a common ordinary collar,
\[
 I_{\rm alg}(\gamma,\omega)=\HInt_\gamma\omega.
\]
\end{proposition}
\begin{proof}
For a finite pole $a=a_0+\delta$, ordinary compactness gives a positive ordinary separation between $\gamma_0$ and $a_0$. The real-coefficient semialgebraic description has the same lower bound on the field interval. Uniform infinitesimal closeness makes the straight homotopy from $(\gamma,a)$ to $(\gamma_0,a_0)$ avoid zero differences throughout. Hence
\[
 \Wind_{\rm sa}(\gamma,a)=\Wind(\gamma_0,a_0).
\]
Equation \eqref{eq:wind} gives equality of the pairings on $dz/(z-a)$. An infinite normalized pole has winding zero, and its reciprocal has a coherent primitive on a sufficiently large bounded normalized disk, so its Hahn period is also zero. Exact rational differentials have zero period by the coefficientwise fundamental theorem on the admitted pullbacks. Partial fractions now prove the assertion, with affine scale factors cancelling as in \eqref{eq:affineintegral}.
\end{proof}

The explicit overlap hypotheses prevent an unwarranted claim that every smooth Hahn-deformed path extends to a semialgebraic map on $[0,1]_{R_\Gamma}$. Most do not. The proposition says the two theories agree when their objects and their homotopies really can be compared.

For the full semialgebraic unit circle, \eqref{eq:algebraicpairing} immediately gives
\[
 I_{\rm alg}\left(S^1,\frac{dz}{z-(1-t)}\right)=2\pi i,\qquad
 I_{\rm alg}\left(S^1,\frac{dz}{z-(1+t)}\right)=0.
\]
These are legitimate values in a case where the unmodified macroscopic Hahn pullback fails. They do not contradict \cref{prop:comparison}, whose leading-separation hypothesis is absent.

\subsection{What this pairing does not determine}
For an open path, the differential $dz/(z-a)$ requires a logarithmic increment, not just a residue and a winding integer. A choice of path continuation or a more elaborate period object is needed. A global fine logarithm is not automatically that choice, as \cref{sec:obstructions} shows. For nonrational analytic functions, partial fractions need not leave a rational exact remainder, so \eqref{eq:algebraicpairing} is not a definition on all of $\HH(U)$ or $\A_n$. Those larger classes are handled by the coherent and microscopic constructions, not by pretending that the rational formula already covers them.

\section{Other research approaches and their exact scope}\label{sec:alternatives}

\subsection{Semialgebraic chains, Stokes, and the target of integration}
Ohmoto--Shiota prove a $C^1$ triangulation theorem and state Stokes for oriented semialgebraic chains with fundamental cycles \cite[Theorem~4.1 and Proposition~4.2]{OS}. This is an important existing counterpart, not merely a proposal. Their non-Archimedean discussion says that integral values may require a larger saturated field and refers to further measure theory. It should not be read as a construction of the particular $K$-valued contour operator defined in this article.

There is a further codomain distinction in the cited literature. Ma\v r\'{i}kov\'{a}--Shiota construct a measure taking values in an ordered semiring, a Dedekind completion of a quotient, and prove a volume-preserving change-of-variables property \cite{MS}. Berarducci--Otero's earlier additive measure is real-valued on definable subsets of the finite part \cite{BOmeasure}. These are different targets and retain different information. A reference to any of them is not a proof that every signed or complex differential form has a canonical integral in the original Hahn field.

For a surcomplex application one must therefore specify the scalar integration operator on each simplex, its target, and an embedding or comparison map if one wants values in $\SC$. Triangulation supplies a finite chain decomposition; it does not by itself solve the scalar integration problem. In particular, saturation alone does not make an ordinary sequence $1/n$ converge to zero in the full ordered topology of a field containing infinitesimals. The finite-mesh obstruction remains unless the limiting notion has been changed.

A promising combination is to use definable triangulation to organize domains and then impose Hahn-coherent pullback conditions on each simplex. Theorems \ref{thm:stokes} and \ref{thm:deformedstokes} give Stokes whenever those checks succeed. A theorem that all desired definable data admit such a finite coherent atlas would be an additional result, not a consequence of triangulation alone.

\subsection{Berkovich geometry changes the space, not just the topology}
For a rank-one subgroup $\Gamma\subset\R$, put $|x|_v=\exp(-v(x))$. The Hahn field is complete and algebraically closed, so it is a legitimate base for Berkovich geometry. The Berkovich line adds multiplicative seminorms on $K[T]$, not just more field points. Its affine line is tree-like and uniquely arc-connected; deleting a branch point can create infinitely many components \cite[\S2.3.3]{Temkin}. It is therefore not a replacement complex plane with the same circle-separation geometry.

Chambert-Loir--Ducros construct real $(p,q)$-forms and currents using tropicalizations and skeleta, with integration of top bidegree forms and a Stokes formula. The revised 2025 version develops this in a broad tropical-space setting \cite{CLD}. This is a powerful established analogue of part of complex differential geometry. Its real-valued forms, added analytic points, and integration degrees differ from the $K$-valued holomorphic $1$-form periods of \cref{sec:cauchy}.

The rank restriction is substantive. An arbitrary higher-rank $\Gamma$ cannot be embedded order-preservingly into $\R$. Passing to a real-valued coarsening may lose infinitesimal scales, while retaining them requires a different higher-rank geometric framework. Also $|c|_v=1$ for every nonzero ordinary complex constant $c$, however small its usual complex modulus. Thus ordinary coefficient convergence is not non-Archimedean convergence. None of these differences prevents Berkovich methods from being useful; they explain why a Berkovich Stokes theorem cannot be quoted as the proof of the specific contour identities in the supplied manuscripts.

\subsection{Roots-of-unity quadrature and the danger of the wrong limit}
Roots-of-unity averages suggest a discrete contour construction, analogous in spirit to the averaging methods used in Shnirelman-type non-Archimedean integration \cite{Ettayb}. Rather than import a $p$-adic convergence theorem, consider the following identity directly. For $f$ on a circle of center $a$ and radius $r$, set
\begin{equation}\label{eq:rootaverage}
 Q_N(f)=\frac1N\sum_{\zeta^N=1}r\zeta f(a+r\zeta),
\end{equation}
where the roots of unity are the ordinary complex ones. If $f(a+w)=\sum_{j=j_0}^{j_1}c_jw^j$ is a Laurent polynomial, then
\[
 Q_N(f)=\sum_{\substack{j_0\le j\le j_1\\N\mid j+1}}c_jr^{j+1}.
\]
This follows from $N^{-1}\sum_{\zeta^N=1}\zeta^q=1$ for $N\mid q$ and $0$ otherwise. When $N$ exceeds every nonzero $|j+1|$ in the finite support, the average equals $c_{-1}$. The corresponding normalized contour integral is exactly that residue.

For ordinary holomorphic coefficient functions on an annulus, the ordinary complex limit of each coefficient average recovers its ordinary contour integral. Taking these ordinary limits \emph{coefficientwise} gives the Hahn contour functional. That is a legitimate alternative description, but it is not a limit in the internal valuation topology.

\begin{example}[A failed valuation-limit contour definition]\label{ex:quadrature}
For $a=0$, $r=1$, and $f(z)=1/(z-b)$, with the ordinary real constant $b=1/2$, logarithmic differentiation of $X^N-1$ gives
\[
 Q_N(f)=\frac1N\sum_{\zeta^N=1}\frac\zeta{\zeta-b}
 =\frac1{1-b^N}.
\]
The ordinary complex limit is $1$, but $v(Q_N(f)-1)=0$ for every $N$. Hence it does not converge to $1$ in the Hahn valuation topology. In the full fine topology it cannot converge at all because it is not eventually constant. If instead $b=t^\rho$, then $v(Q_N(f)-1)=N\rho$: a rank-one valuation can see convergence, but a higher-rank group with an element larger than every $N\rho$ still cannot.
\end{example}

This elementary calculation isolates the obstruction to an automatic Shnirelman transplant. The coefficient and valuation geometries must be treated separately. Finite Laurent data give exact algebraic quadrature; general coherent coefficients give ordinary coefficient limits; valuation convergence needs an additional cofinality and regularity argument.

\subsection{Surreal antiderivatives and resurgent extensions}
Costin--Ehrlich construct extensions and integration for substantial classes of resurgent real functions, with treatments of singularities and ordered exponential subfields \cite{CE}. Their framework is a natural candidate when a contour is parameterized by functions whose real and imaginary pullbacks belong to those classes. It can retain asymptotic information that a mere standard-part integral would lose.

To turn that idea into a contour theorem, however, one must prove closure under the required compositions $F(\gamma(s))\gamma'(s)$, specify logarithmic branches and endpoint conventions, and establish a substitution and homotopy theorem. Agreement with Hahn integration should first be tested on common analytic germs, rational forms, and the microscopic residues above. The real-variable integration results alone do not establish a Jordan theorem or a global complex contour theory.

There is no conflict between the broadness of some such function-extension theories and the arbitrary smooth coefficient forms used on ordinary parameter manifolds in \cref{sec:stokes}. The latter construction integrates ordinary functions first and stores their values as Hahn coefficients. It does not assert a canonical fine-smooth extension of every ordinary smooth function to every surcomplex point.

\subsection{Comparison of the integration strategies}
\begin{center}\small
\begin{tabular}{@{}L{0.22\textwidth}L{0.33\textwidth}L{0.35\textwidth}@{}}
\toprule
Strategy & What it supplies & What must not be assumed\\
\midrule
Hahn coefficient integration & Exact $K$-valued periods and Stokes on admitted parameter chains & Fine continuity or arbitrary pole avoidance is sufficient\\[0.35em]
Microscopic rescaling & Contour realizations for radius-free analytic-germ residues & Original coefficient germs share an ordinary radius\\[0.35em]
Algebraic winding and residues & Rational closed-contour pairing at arbitrary scales & A prior Riemann integral or canonical open-path logarithms\\[0.35em]
Semialgebraic integration & Finite triangulated chains and Stokes formulations & Every version has the same scalar target\\[0.35em]
Berkovich--tropical integration & Real differential forms on enlarged analytic spaces & The original point plane or all higher valuation ranks are preserved\\[0.35em]
Surreal antiderivatives & Integration on distinguished extendable function classes & Automatic closure under arbitrary contour substitutions\\[0.35em]
Roots-of-unity averaging & Exact Laurent extraction; coefficientwise quadrature & Ordinary coefficient limits are valuation limits\\
\bottomrule
\end{tabular}
\end{center}

\section{Worked calculations at several scales}\label{sec:examples}

\subsection{An infinitesimal rectangle has nonzero area}
Let $A,B\in R_\Gamma$ be positive, and parameterize a rectangle with side lengths $A,B$ by the ordinary unit square. For $\alpha=-y\dd x+x\dd y$,
\[
 d\alpha=2\dd x\wedge\dd y.
\]
The right edge contributes $AB$, the top edge contributes $AB$, and the other edges contribute zero. Therefore
\begin{equation}\label{eq:rectangle}
 \HInt_{\partial([0,A]\times[0,B])}\alpha
 =2AB
 =\HInt_{[0,A]\times[0,B]}d\alpha.
\end{equation}
The chain notation denotes the indicated affine parameterization; it does not presume a fine-continuous ordinary square. For $A=t$ and $B=t^\omega$, the value is $2t^{1+\omega}\ne0$. Standard part alone would make it zero, losing the exact Stokes identity's geometric content.

\subsection{A deformed ellipse: holomorphic and nonholomorphic forms differ}
Let $R=1+t$ and $S=1+t^\omega$. On the ordinary disk use the deformation $(x,y)\mapsto(Rx,Sy)$, with boundary
\[
 \gamma(\theta)=R\cos\theta+iS\sin\theta.
\]
Direct multiplication gives
\[
 \overline\gamma\gamma'
 =(S^2-R^2)\sin\theta\cos\theta+iRS.
\]
Thus
\begin{equation}\label{eq:ellipse}
 \HCoint_\gamma\overline z\dd z=2\pi i RS
 =2\pi i(1+t+t^\omega+t^{1+\omega}).
\end{equation}
Stokes gives the same answer from $d(\overline z\dd z)=2i\dd x\wedge\dd y$. In contrast, $z^q\dd z$ is exact for every integer $q\ge0$, and its integral is zero for both the circle and the ellipse. This example tests the distinction between general Stokes and closed-form deformation invariance.

\subsection{Two infinite residues with a finite sum}
For $F(z)=e^z/(z^2-t)$ the two simple residues are
\[
 \frac{e^{t^{1/2}}}{2t^{1/2}},\qquad
 -\frac{e^{-t^{1/2}}}{2t^{1/2}}.
\]
They are individually infinite, but their sum is
\begin{equation}\label{eq:expresidue}
 \frac{\sinh(t^{1/2})}{t^{1/2}}
 =\sum_{n\ge0}\frac{t^n}{(2n+1)!}.
\end{equation}
The ordinary unit circle and the protecting microscopic circle $|z|=t^{1/4}$ both have integral $2\pi i$ times \eqref{eq:expresidue}. The individual microscopic circles described in \cref{sec:scales} recover the separate weights. This calculation is the moving-pole example from \cite{Analysis}, now also illustrating the rescaling bridge.

\subsection{A radius-free numerator with a computable contour period}\label{subsec:radiusfreeexample}
Take the numerator \eqref{eq:shrinking} and denominator $F(z)=z^2-t$. For the protecting radius $r=t^{1/4}$, \cref{thm:microcircle} gives
\begin{align}\label{eq:radiusfreeexample}
 \frac1{2\pi i}\HCoint_{rS^1}\frac{H(z)}{z^2-t}\dd z
 &=\sum_{k\ge0}t^k[z^{2k+1}]H(z)\notag\\
 &=\sum_{m\ge1}\sum_{k\ge0}m^{2k+1}t^{m+k}
 =\sum_{m\ge1}\frac{m t^m}{1-m^2t}.
\end{align}
Every rearrangement is strongly valid: the coefficient at $t^N$ has only the pairs $m+k=N$, $m\ge1$, $k\ge0$. The first terms are
\begin{equation}\label{eq:radiusfreenumbers}
 t+3t^2+12t^3+64t^4+441t^5+3855t^6+41464t^7+533736t^8+\cdots.
\end{equation}
Equivalently, the actual simple-pole formula is
\[
 \frac{H(t^{1/2})-H(-t^{1/2})}{2t^{1/2}}.
\]
For each summand of $H$, this difference quotient is $m t^m/(1-m^2t)$, proving agreement with \eqref{eq:radiusfreeexample}. No common ordinary disk for the original numerator coefficients exists. The period is therefore genuinely supplied by the microscopic representation, not by an unavailable original fixed-circle integral.

\subsection{A two-variable torus and cancellation across valuation ranks}
Consider the source example
\[
 F_1=x^2-Ay,\qquad F_2=y^2-Bx,\qquad A,B\in\m\setminus\{0\}.
\]
Choose $2\rho_x<v(A)$ and $2\rho_y<v(B)$. The torus formula gives
\begin{equation}\label{eq:coupledres}
 \mathcal R_F(H)=\sum_{k,\ell\ge0}A^kB^\ell
 [x^{2k+1-\ell}y^{2\ell+1-k}]H,
\end{equation}
where a coefficient with a negative exponent is zero. In particular,
\[
 \mathcal R_F(1)=0,\quad \mathcal R_F(xy)=1,\quad
 \mathcal R_F(4xy-AB)=4.
\]
The Jacobian is $J_F=4xy-AB$. The origin has Jacobian $-AB$. The other three roots have $xy=AB$ and Jacobian $3AB$. Their weights for numerator $1$ are therefore
\[
 -\frac1{AB},\quad\frac1{3AB},\quad\frac1{3AB},\quad\frac1{3AB},
\]
which sum to zero. For $A=t$ and $B=t^\omega$, this cancellation involves weights of valuation $-(1+\omega)$, whereas the torus result has no negative term at all. Equation \eqref{eq:coupledres} is the residue formula from \cite{Geometry}; \cref{thm:microtorus} is its explicit microscopic contour realization.

\subsection{What the computational checks verify}
The accompanying script \texttt{verify\_examples.py} uses exact SymPy arithmetic to check the rectangle and ellipse calculations, the first eight coefficients in \eqref{eq:radiusfreenumbers}, the residue coefficient formula for the coupled system, its trace--Jacobian identities on a finite monomial test set, finite roots-of-unity selection, and the finite-partition square identity. The recorded run passes $317$ exact checks. The radius-free coefficient list was also independently checked with a Wolfram Language evaluation. These checks are tests of displayed finite identities and truncations, not implementations of arbitrary Hahn fields and not formal verification of the general theorems.

\section{A practical framework and remaining research questions}\label{sec:outlook}

\subsection{How to select a contour theory for a concrete problem}
For a common-domain Hahn function, first choose an ordinary base contour whose reduction avoids the leading singularities. The coefficientwise integral and \cref{thm:displacement} then handle all admitted infinitesimal changes of shape and endpoints. For a radius-free germ near an isolated zero cluster, first choose positive protecting scales and apply \cref{thm:microcircle} or \cref{thm:generalrepresentation}. For rational functions on full semialgebraic loops, use \eqref{eq:algebraicpairing} when an ordinary collar is unavailable. If the purpose is separation rather than integration, specify whether the desired sides are coarse residue sides or definable sides in the real closed plane.

At each step, record the coefficient algebra, exponent group, parameter domain, differential, and integral target. Those data remove most apparent paradoxes. A proof that uses a derivative-zero fine function, an ordinary real contour, a definable Jordan interior, and a valuation limit without comparison theorems is generally mixing four different constructions.

\subsection{Research problems sharpened by the results}
A natural next problem is a \emph{boundary-layer atlas theorem}: under explicit tame hypotheses, can rational or selected analytic data with poles infinitesimally close to a leading contour be covered by finitely many scale charts whose pullback integrals glue? The values for rational closed contours are already fixed by \eqref{eq:algebraicpairing}; the issue is realizing them by a uniform chain construction which also accommodates open paths and nonrational forms.

A second problem is a \emph{comparison theorem with definable integration}. One should identify a common class on which a chosen semialgebraic integration operator, with its actual codomain specified, maps naturally to the Hahn functional. Polynomial simplex forms are a necessary initial test. Logarithmic forms and infinitesimal boundary layers are more discriminating tests because choices or loss of scale can become visible.

A third is an \emph{intrinsic coherent chain category}. The present construction uses explicit ordinary parameterizations and support-controlled pullbacks. A useful global theory would define admissible atlases and equivalence of chains, prove a satisfactory gluing theorem, and identify their homology independently of a chosen presentation. The proof must address possibly unbounded collections of supports rather than assuming that local coherence gives global common-support data.

A fourth is a \emph{higher-rank analytic comparison}. The microscopic torus formulas retain every exponent in $\Gamma$. A relationship with analytic skeleta or higher-rank valuations should preserve those exact periods, or state precisely which quotient of the scale information it keeps. Merely applying a rank-one absolute value cannot provide that comparison.

These are research directions formulated by the present analysis, not statements that a particular named published conjecture remains open. The literature search was conducted on September 21, 2026, and was targeted rather than exhaustive. We make no assertion that the exact support-controlled package here has an established priority over all analytic-structure or generalized-series frameworks.

\subsection{Conclusion}
There is no single unqualified Jordan--Cauchy--Stokes theorem on the raw surcomplex point class. There are, however, precise and useful counterparts. Definable topology restores finite-dimensional separation and integer winding. Hahn-coherent forms and chains restore exact periods and generalized Stokes without relying on fine-continuous paths. Positive microscopic rescaling makes even radius-free analytic germs accessible to contour methods, and residue transformation represents arbitrary isolated complete-intersection residues by controlled microscopic tori. Rational winding pairings fill an important part of the gap where macroscopic coefficientwise pullbacks fail.

The unifying principle is not to discard infinitesimals or force every scale into one real-valued metric. It is to specify which information each topology retains and which finite-at-each-exponent operations justify each integral identity.

\appendix
\section{Support certificates and logical dependencies}\label{app:support}

The following table records bounds used in the proofs. $S$ is a source support, $E$ a positive error support, and $P=\{\rho_1,\ldots,\rho_n\}$ a finite positive rescaling alphabet.
\begin{center}\small
\begin{tabular}{@{}L{0.57\textwidth}L{0.34\textwidth}@{}}
\toprule
Operation & Containing support\\
\midrule
Ordinary linear integration, differentiation, residue, or a chosen linear primitive operator & $S$\\[0.4em]
Product of two Hahn families & $S+T$\\[0.4em]
Infinitesimal Taylor evaluation or pullback & $S+E^*$\\[0.4em]
Positive monomial rescaling & $S+P^*$\\[0.4em]
Perturbation residue after integration & $S+E^*$\\[0.4em]
Normalized coordinate-power denominator error in component $j$ & $(\supp E_j-d_j\rho_j)+P^*$\\
\bottomrule
\end{tabular}
\end{center}
The last containing set is positive under the protecting-radius inequality. Before torus integration the numerator and the volume factors may introduce a fixed shift, but a shift of a well-ordered support is still well ordered. After integration the scale factors cancel, giving the simpler residue support stated in the table.

For clarity, ``only finitely many dependencies at an exponent'' does not mean ``only finitely many exponents below it.'' The latter fails, for example, below $\omega$ in $\Q+\Q\omega$. The proofs use the former statement, which follows from finite word decompositions. This distinction is essential in the residue transformation argument and in the passage from infinitely many germs to one finite ordinary-parameter calculation.

The dependency structure is also explicit. The support lemma defines the Hahn operations. Ordinary Stokes proves Hahn Stokes. Taylor pullback and finite-jet identities prove the deformed-chain version. Cauchy follows from closedness and ordinary coefficient Cauchy. Independently, positive rescaling and Laurent coefficient extraction realize the perturbation residue. Classical finite-family residue transformation extends the realization to general isolated complete intersections. The rational pairing is defined independently from semialgebraic winding and is compared with Hahn integration only afterward. No Jordan statement in the fine topology and no unproved general fine integration rule enters this chain of arguments.

\section{Source audit and reproducibility}\label{app:audit}
The two supplied manuscripts are treated as research manuscripts, not as refereed references. The present construction uses their precise definitions of $\HH(U)$, $\A_n$, strong evaluation, and the perturbation residue. We provide proofs of the added chain, homotopy, and microscopic representation statements. Claims about all maximal ideals, regularity, faithful flatness, and complete zero-scheme conservation in the geometry manuscript are not silently used to establish Stokes; the optional finite-algebra interpretation is explicitly attributed in \cref{eq:traceresidue}.

The closest external inputs are the real-closed-field winding and separation literature, ordinary differential topology and residue theory, and existing non-Archimedean integration frameworks. In particular, the semialgebraic Stokes reference does not eliminate the need to state a target for the integral; the Berkovich reference does not identify its spaces with the original point plane; and the surreal antiderivative reference does not provide arbitrary complex contour substitution for free.

The archive contains this \LaTeX\ source, its compiled PDF, an exact-arithmetic verification script and recorded output, a small Wolfram Language cross-check, and build instructions. It does not contain the external publications or redistribute font files. The supplied source filenames are listed in the bibliography so that the direction of dependence can be checked against the original attachments.

\begin{thebibliography}{99}
\addcontentsline{toc}{section}{References}

\bibitem{Analysis}
\emph{Surcomplex Analysis: Infinitesimal Calculus, Hahn-Coherent Holomorphy, and Contour Theory}.
Supplied research manuscript, September 21, 2026.
Source file: \texttt{surcomplex\_analysis(3).tex}.
Especially the sections on fine topology, coherent contour functionals, moving-pole residues, global fine logarithms, and affine scales.

\bibitem{Geometry}
\emph{Infinitesimal Analytic Geometry over the Surcomplex Numbers: Radius-Free Hahn Germs, a Nullstellensatz, and Conservation of Singular Zeros}.
Supplied research manuscript, September 21, 2026.
Source file: \texttt{surcomplex\_analytic\_geometry.tex}.
Especially the definitions of $\A_n$, strongly summable perturbation residues, and finite dependency at a fixed exponent.

\bibitem{BCR}
J. Bochnak, M. Coste, and M.-F. Roy,
\emph{Real Algebraic Geometry}, Ergebnisse der Mathematik und ihrer Grenzgebiete (3), vol.~36, Springer, 1998.
\href{https://doi.org/10.1007/978-3-662-03718-8}{doi:10.1007/978-3-662-03718-8}.
Semialgebraic triangulation and topology over real closed fields.

\bibitem{BOintersection}
A. Berarducci and M. Otero,
``Intersection theory for o-minimal manifolds,''
\emph{Annals of Pure and Applied Logic} \textbf{107} (2001), 87--119.
\href{https://doi.org/10.1016/S0168-0072(00)00027-0}{doi:10.1016/S0168-0072(00)00027-0}.
\href{https://arpi.unipi.it/handle/11568/186193}{Institutional record}.

\bibitem{BOhomology}
A. Berarducci and M. Otero,
``O-minimal fundamental group, homology and manifolds,''
\emph{Journal of the London Mathematical Society} (2) \textbf{65} (2002), 257--270.
\href{https://doi.org/10.1112/S0024610701003015}{doi:10.1112/S0024610701003015}.

\bibitem{BOmeasure}
A. Berarducci and M. Otero,
``An additive measure in o-minimal expansions of fields,''
\emph{Quarterly Journal of Mathematics} \textbf{55} (2004), 411--419.
\href{https://doi.org/10.1093/qmath/hah010}{doi:10.1093/qmath/hah010}.

\bibitem{CDS}
E. Cattani, A. Dickenstein, and B. Sturmfels,
``Computing multidimensional residues,'' in \emph{Algorithms in Algebraic Geometry and Applications}, Progress in Mathematics, vol.~143, Birkh\"auser, 1996, 135--164.
\href{https://arxiv.org/abs/alg-geom/9404011}{arXiv:alg-geom/9404011};
\href{https://doi.org/10.1007/978-3-0348-9104-2_8}{doi:10.1007/978-3-0348-9104-2\_8}.
The introduction records local residues and the transformation law; later sections discuss normal forms and traces.

\bibitem{CE}
O. Costin and P. Ehrlich,
``Integration on the surreals,''
\emph{Advances in Mathematics} \textbf{452} (2024), 109823.
\href{https://doi.org/10.1016/j.aim.2024.109823}{doi:10.1016/j.aim.2024.109823};
\href{https://arxiv.org/abs/2208.14331}{arXiv:2208.14331}, version~5, July 4, 2024.

\bibitem{CLD}
A. Chambert-Loir and A. Ducros,
\emph{Formes diff\'erentielles r\'eelles et courants sur les espaces de Berkovich},
\href{https://arxiv.org/abs/1204.6277}{arXiv:1204.6277}, version~2, July 25, 2025.
The revised manuscript develops real forms, tropical spaces, integration, and a Stokes formula.

\bibitem{Eisermann}
M. Eisermann,
``The fundamental theorem of algebra made effective: an elementary real-algebraic proof via Sturm chains,''
\href{https://arxiv.org/abs/0808.0097}{arXiv:0808.0097}, version~4, March 26, 2012.
Especially Theorem~1.2 on algebraic winding and Theorem~1.5 on root counting.

\bibitem{Ettayb}
J. Ettayb,
``Some integrals for groups of bounded linear operators on finite-dimensional non-Archimedean Banach spaces,''
\emph{Buletinul Academiei de \c Stiin\c te a Republicii Moldova. Matematica} (2022), no.~3, 3--14.
\href{https://doi.org/10.56415/basm.y2022.i3.p3}{doi:10.56415/basm.y2022.i3.p3}.
A reference for contemporary use of Volkenborn and Shnirelman integration; no convergence theorem from it is imported into the Hahn construction.

\bibitem{Lee}
J. M. Lee,
\emph{Introduction to Smooth Manifolds}, 2nd ed., Graduate Texts in Mathematics, vol.~218, Springer, 2013.
\href{https://doi.org/10.1007/978-1-4419-9982-5}{doi:10.1007/978-1-4419-9982-5}.
Ordinary exterior calculus, Stokes, and de Rham theory.

\bibitem{MS}
J. Ma\v r\'{i}kov\'{a} and M. Shiota,
``Measuring definable sets in o-minimal fields,''
\emph{Israel Journal of Mathematics} \textbf{209} (2015), 687--714.
\href{https://doi.org/10.1007/s11856-015-1234-0}{doi:10.1007/s11856-015-1234-0};
\href{https://arxiv.org/abs/1402.4787}{arXiv:1402.4787}.
The measure target is an ordered semiring, not the original scalar field.

\bibitem{Neumann}
B. H. Neumann,
``On ordered division rings,''
\emph{Transactions of the American Mathematical Society} \textbf{66} (1949), 202--252.
The ordered-series support lemma is the algebraic summability input.

\bibitem{OS}
T. Ohmoto and M. Shiota,
``$C^1$-triangulations of semialgebraic sets,''
\emph{Journal of Topology} \textbf{10} (2017), 765--775.
\href{https://doi.org/10.1112/topo.12024}{doi:10.1112/topo.12024};
\href{https://arxiv.org/abs/1505.03970}{arXiv:1505.03970}, version~2, August 7, 2017.
See Section~4 for triangulated integration and Stokes, and the introduction for the non-Archimedean codomain qualification.

\bibitem{Poonen}
B. Poonen,
``Maximally complete fields,''
\emph{L'Enseignement Math\'ematique} (2) \textbf{39} (1993), 87--106.
\href{https://math.mit.edu/~poonen/papers/amsval.pdf}{Author's full text}.
See the Hahn-field construction and Corollary~4.

\bibitem{Temkin}
M. Temkin,
``Introduction to Berkovich analytic spaces,''
\href{https://arxiv.org/abs/1010.2235}{arXiv:1010.2235}, version~2, December 16, 2011.
Especially Section~2.3.3 on the analytic affine line and its tree structure.

\bibitem{Tsikh}
A. K. Tsikh,
\emph{Multidimensional Residues and Their Applications},
Translations of Mathematical Monographs, vol.~103, American Mathematical Society, 1992.
\href{https://doi.org/10.1090/mmono/103}{doi:10.1090/mmono/103}.

\bibitem{vdDE}
L. van den Dries and P. Ehrlich,
``Fields of surreal numbers and exponentiation,''
\emph{Fundamenta Mathematicae} \textbf{167} (2001), 173--188;
erratum, \textbf{168} (2001), 295--297.
\href{https://doi.org/10.4064/fm167-2-3}{doi:10.4064/fm167-2-3}.

\end{thebibliography}
\end{document}
